% Propietario conservado: /Users/ruben/Documents/New project/output/REVISION_ARGUMENTAL_APERTURAS_CIERRES_20260915/IX/ES/sections/ch_cierre_cardinal.tex
\section{Démonstration de \texorpdfstring{\(2^{\aleph_0}=\aleph_1\)}{2 élevé à aleph zéro égal à aleph un} dans le continu HMT}
\label{ix:sec:cierre-cardinal-nativo}

Cette section réunit en une démonstration contiguë le prolongement, la double projection et la profondeur arbitraire ; elle construit ensuite la clôture sémantique, le relèvement uniforme, la machinerie généalogique et la clôture cardinale. La hiérarchie ensembliste n'est identifiée qu'au terme de la construction et ne peut être lue comme une origine causale. La dépendance causale est :

\[
 \mathrm{APP}\longrightarrow\mathrm{TRIT}\longrightarrow\mathrm{TPK}
 \longrightarrow X_\infty^{\rm enr}
 \longrightarrow\mathcal C_{\rm cont}^{\rm disc}
 \longrightarrow m_\infty
 \longrightarrow\mathfrak G_{\rm HMT}(h,m_\infty).
\tag{32}\label{ix:eq:ch-cadena-causal}
\]

On construit d’abord l’état, ses prolongements, le continu et le registre généalogique.
Ensuite, on prolonge \textbf{cette généalogie produite} au moyen d’une règle sémantique
explicite à tous ses objets héréditaires. L’identification ultérieure à
une hiérarchie ensembliste connue ne sélectionne ni les germes, ni les opérateurs,
ni le registre généalogique, ni la droite, ni ses projections. Elle permet bien de vérifier, en toute
transparence, la force logico-cardinale exacte de la règle sémantique adoptée.

\subsection{Prolongement d’états et extension sémantique transfinie}

Soit

\[
 \Sigma_{\rm HMT}
 =\bigl(\dot R_{\rm APP},\dot R_{\rm TRIT},\dot R_{\rm TPK},
 \dot R_{U},\dot R_{\Gamma_9},\dot R_{\rm Ext},\dot R_\rho,
 \dot R_{\pi},\dot R_{\varphi},\dot R_e,\dot R_{\alpha},
 \dot R_{\rm ar},\dot R_{\rm exc}\bigr)
\tag{33}\label{ix:eq:ch-signatura}
\]

la signature \textbf{relationnelle} dénombrable dont le code effectif est \(h\). Chaque
\(\dot R_S\) est le graphe de l’opérateur \(S\) sur les codes où il est
défini ; on ne présuppose pas qu’une fonction externe soit totale à l’intérieur de chaque
niveau. Les graphes \(\dot R_{\pi},\dot R_{\varphi},\dot R_e,\dot R_{\alpha}\)
sont les lecteurs de sortie construits par la génération coinductive et la
fermeture dodécaphasique ; ils ne reçoivent pas comme arguments les valeurs conventionnelles des
constantes. Leur présence dans \(\Sigma_{\rm HMT}\) assure que la clôture
conserve l’état complet : elle intègre simultanément les routes qui publient la
valeur et l’incidence qui se prolonge vers la réalisation exceptionnelle.

À chaque horizon, le code de l'état enregistre les coordonnées APP--TRIT--TPK, le chemin, le feuillet, la retenue, la frontière et la mémoire ; à ses étapes de publication, il enregistre également \(P_N,\Phi_N,E_N,A_N\), où \(A_N\) est le préfixe de \(\alpha_{\rm HMT}\) produit par le même état. Ces coordonnées ne sélectionnent pas la branche de l'extérieur : elles font partie de l'histoire qui est codée. Une numérotation primitive récursive des coordonnées finies d'un état étant fixée, toute histoire inverse \(x_\infty=(x_N)_{N\geq1}\in X_\infty^{\rm enr}\) détermine le code
\[
 \operatorname{code}(x_\infty)
 :=\{\langle N,j,c\rangle:(x_N)_j\text{ a pour code }c\}\subseteq\omega.
\]
Les équations de cohérence
\(\rho_{N+1,N}(x_{N+1})=x_N\) permettent de décoder de manière unique
l’histoire à partir de ce réel. Nous écrivons \(m_\infty\) pour le code de
l’histoire réalisée et \(m_\infty(n)\) pour sa fonction caractéristique sous la
numérotation fixée. Le code opératoire \(h\) et ce registre complet sont
réunis, au moyen d’une fonction primitive récursive d’appariement, dans le réel

\[
 u_{h,m}:=
 \{\langle0,\langle n,h(n)\rangle\rangle:n<\omega\}
 \cup
 \{\langle1,\langle n,m_\infty(n)\rangle\rangle:n<\omega\}.
\tag{33a}\label{ix:eq:ch-codigo-conjunto}
\]

Les deux projections primitives récupèrent \(h\) et \(m_\infty\) à partir de \(u_{h,m}\).

La première extension est le prolongement métrique et opératoriel des états finis. Pour \(N\ge1\), soit \(\widehat{\mathcal X}^{\rm enr,raw}_{6N}\) l'ensemble des états enrichis bruts de longueur \(6N\), et définissons en termes ensemblistes le sous-système survivant par
\[
\begin{aligned}
 X_N
 &:=\widehat{\mathcal X}^{\rm enr,surv}_{6N}\\
 &:=\Bigl\{x\in\widehat{\mathcal X}^{\rm enr,raw}_{6N}:\\[-.2ex]
 &\qquad\forall M>N\ \exists z\in
   \widehat{\mathcal X}^{\rm enr,raw}_{6M}\;\text{ tel que}\\[-.2ex]
 &\qquad\operatorname{tr}_{6M,6N}(z)=x\ \wedge\
   z\text{ satisfait toute transition TPK intermédiaire}\Bigr\}.
\end{aligned}
\]
et posons
\[
 \rho_{N+1,N}:=\operatorname{tr}_{6N+6,6N}\!\upharpoonright X_{N+1},
\]
et
\[
 \operatorname{Ext}_N(x):=
 \{y\in X_{N+1}:(x,y)\in
 \operatorname{Ext}_{6N+6,6N}\}.
\]
Le niveau racine est inclus sans exception implicite :
\[
 X_0:=\{\ast\},\qquad
 \rho_{1,0}:X_1\to X_0\text{ est constant},\qquad
 \operatorname{Ext}_0(\ast):=X_1.
\]
L’exposant \(\mathrm{surv}\) désigne
l’arbre élagué des états qui possèdent des descendants TPK compatibles à tout
horizon fini. La loi de prolongement TPK établie dans les sections
précédentes prouve que les germes
émis appartiennent à cet arbre ; la ramification finie et le lemme de Kőnig
transforment la compatibilité à tout horizon en une branche infinie. En
particulier,

\[
 \forall N\ge0\;\forall x\in X_N\quad
 \varnothing\ne\operatorname{Ext}_N(x)
 \subseteq\rho_{N+1,N}^{-1}(x).
\tag{34}\label{ix:eq:ch-extension-no-vacia}
\]

Par le même lemme de l’arbre à ramification finie, tout germe survivant
possède un prolongement compatible

\[
 x_\infty=(x_N)_{N\ge1}\in
 X_\infty^{\rm enr}:=\varprojlim_N(X_N,\rho_{N+1,N}).
\tag{35}\label{ix:eq:ch-limite-inverso}
\]

La seconde extension agit \textbf{ensuite} sur le code complet de cette
histoire réalisée.
Elle ne s’infère pas de la seule existence de l’opérateur fini
\(\operatorname{Ext}_N\) : elle formalise à l’échelle ensembliste le principe HMT selon
lequel seuls sont objets du domaine intégral ceux qui possèdent une dérivation
généalogique depuis l’état et ses règles. L’extension sémantique transfinie
est :

\[
 \begin{aligned}
  \mathfrak G_0(h,m_\infty)
  &:=\operatorname{TC}\bigl(\{\omega,u_{h,m},h,m_\infty\}\bigr),\\
  \mathfrak G_{\beta+1}(h,m_\infty)
  &:=\operatorname{Def}_{\Sigma_{\rm HMT}}
       \bigl(\mathfrak G_\beta(h,m_\infty)\bigr),\\
  \mathfrak G_\lambda(h,m_\infty)
  &:=\bigcup_{\beta<\lambda}\mathfrak G_\beta(h,m_\infty)
       \quad(\lambda\text{ limite}),\\
  \mathfrak G_{\rm HMT}(h,m_\infty)
  &:=\bigcup_{\beta\in\operatorname{Ord}}
       \mathfrak G_\beta(h,m_\infty).
 \end{aligned}
\tag{36}\label{ix:eq:ch-clausura-transfinita}
\]

L'hypothèse du continu et une bijection cible ne sont pas introduites parmi les conditions qui définissent l'opérateur successeur ; l'opérateur est donné prospectivement par

\[
 \operatorname{Def}_{\Sigma_{\rm HMT}}(M)
 :=M\cup
 \Bigl\{
   \{x\in M:\langle M,\in,
       (\dot R_S\cap M^{k_S})_{S\in\Sigma_{\rm HMT}}\rangle
       \models\varphi(x,\bar p)\}:
   \ulcorner\varphi\urcorner\in\omega,
   \ \bar p\in M^{<\omega}
 \Bigr\}.
\tag{37}\label{ix:eq:ch-operador-def}
\]

Autrement dit : il publie simultanément tous les sous-ensembles définissables par une formule finie du langage relationnel HMT et des paramètres que la généalogie a déjà produits. Il ne consulte ni valeur réelle cible, ni constante conventionnelle, ni cardinal cible, ni modèle extérieur. La formule~\eqref{ix:eq:ch-operador-def} définit la clôture transfinie sur les germes, les opérateurs, l'état enrichi, le registre et la double projection construits précédemment.

\begin{lemma}[Codage conjoint du programme HMT et du registre réalisé]
\label{ix:thm:ch-6}
La paire formée par
le code opératoire \(h\) et le registre \(m_\infty\) est codée par un unique
réel \(u_{h,m}\subseteq\omega\), et toute formule de \(\Sigma_{\rm HMT}\) se
traduit uniformément en une formule du langage \(\{\in,u_{h,m}\}\).

\end{lemma}
\begin{proof}
Tant \(h\) que \(m_\infty\) sont des suites de naturels : \(h\)
encode la signature, les tables finies et les programmes des opérateurs ;
\(m_\infty\) encode une seule histoire réalisée, non l’espace complet des
histoires. La formule~\eqref{ix:eq:ch-codigo-conjunto} et ses deux décodeurs prouvent la première
affirmation. Chaque symbole de~\eqref{ix:eq:ch-signatura} désigne sa \textbf{relation de graphe encodée dans \(h\)}, jamais le graphe externe de toutes ses évaluations possibles. APP,
TRIT, \(U_t\), \(\Gamma_9\), les restrictions \(\rho\) et les extensions
finies ont des graphes primitifs récursifs sur des codes entiers et donc des
formules \(\Delta_0\) absolues entre niveaux transitifs. Les deux
publications s’expriment comme relations de graphe au moyen des formules

\[
 \begin{aligned}
 \dot R_{\rm ar}(x,y)
 &\;\Longleftrightarrow\;
 \forall k\in\omega\;\exists N\;\forall n\ge N\quad
 |\mu_n(x)-y|<2^{-k},\\
 \dot R_{\rm exc}(x,c,z)
 &\;\Longleftrightarrow\;
 \forall N\in\omega\quad z\!\upharpoonright N
   =E_N(x\!\upharpoonright N,c\!\upharpoonright N),
 \end{aligned}
\tag{37c$_0$}\label{ix:eq:ch-grafos-publicacion}
\]

où \(\mu_n\) et \(E_N\) sont les fonctions rationnelles/finies dont les codes
figurent dans \(h\). On remplace par récurrence chaque formule atomique contenant
un symbole HMT par sa formule de graphe de paramètre \(u_{h,m}\). Les
connecteurs et les quantificateurs sont conservés. Cette récurrence sur la complexité
syntaxique produit une traduction \(\varphi\mapsto\varphi^*\) telle que, pour
tout niveau transitif \(M\) de~\eqref{ix:eq:ch-clausura-transfinita} contenant les paramètres,

\[
 \langle M,\in,(\dot R_S\cap M^{k_S})_S\rangle
   \models\varphi(\bar a)
 \quad\Longleftrightarrow\quad
 M\models\varphi^*(\bar a;u_{h,m}).
\tag{37c}\label{ix:eq:ch-eliminacion-uniforme}
\]

Ainsi, \eqref{ix:eq:ch-operador-def} ne peut consulter ni oracle, ni table décimale cible, ni relation non codée introduisant une infinité non dénombrable d'objets en une seule étape.
\end{proof}

La sémantique HMT est maintenant définie \textbf{indépendamment de son identification à \(\mathfrak G\)}. Soit \(\mathcal T_0^{\rm HMT}\) l'ensemble des noms canoniques des éléments de \(\mathfrak G_0(h,m_\infty)=
\operatorname{TC}(\{\omega,u_{h,m},h,m_\infty\})\). En particulier, le niveau initial contient le registre réalisé \(m_\infty\), et non toutes les histoires de \(X_\infty^{\rm enr}\) comme paramètres primitifs. Les termes admissibles sont les arbres bien fondés obtenus par

\[
 \begin{aligned}
 \mathcal T^{\rm HMT}_{\beta+1}
 &:={\mathcal T}^{\rm HMT}_\beta\cup
 \{\langle\beta,\ulcorner\varphi\urcorner,
       t_0,\ldots,t_{k-1}\rangle:
       t_i\in\mathcal T^{\rm HMT}_\beta\},\\
 \mathcal T^{\rm HMT}_\lambda&:=
       \bigcup_{\beta<\lambda}\mathcal T^{\rm HMT}_\beta,
 \qquad
 \mathcal T^{\rm HMT}:=
       \bigcup_{\beta\in\operatorname{Ord}}\mathcal T^{\rm HMT}_\beta.
 \end{aligned}
\tag{37a}\label{ix:eq:ch-terminos}
\]

La valeur d’un terme successeur est le sous-ensemble du niveau précédent défini
par \(\varphi\) dans la structure relationnelle induite, avec les valeurs de
\(t_0,\ldots,t_{k-1}\) comme paramètres. On définit, avant toute
comparaison ensembliste,

\[
 \operatorname{Ob}_{\rm sem}^{\rm HMT}(h,m_\infty)
 :=\{\llbracket t\rrbracket:t\in\mathcal T^{\rm HMT}\}.
\tag{37b}\label{ix:eq:ch-objetos-semanticos}
\]

Telle est la règle d’\textbf{exhaustivité générative} adoptée par HMT : tout
objet du domaine intégral doit posséder l’un de ces arbres de dérivation et tout
arbre admissible publie un objet. Ce n’est pas une conséquence de la ramification
finie. L’hypothèse du continu n’est pas introduite dans cette règle ; elle s’obtient
au moyen de la démonstration développée dans l’article. On n’introduit pas non plus
comme données une énumération des réels ni une bijection cible.

\begin{theorem}[Équivalence entre la sémantique des termes et la hiérarchie généalogique]
\label{ix:thm:ch-6a}
L’évaluation des
termes précédents vérifie

\[
 \operatorname{Ob}_{\rm sem}^{\rm HMT}(h,m_\infty)
 =\mathfrak G_{\rm HMT}(h,m_\infty).
\tag{37d}\label{ix:eq:ch-equivalencia-semantica}
\]

\end{theorem}
\begin{proof}
Par récurrence transfinie. Au niveau initial, les deux classes coïncident. Si toute valeur d'un terme de \(\mathcal T^{\rm HMT}_\beta\) appartient à \(\mathfrak G_\beta\), la sémantique du terme successeur est l'un des sous-ensembles de~\eqref{ix:eq:ch-operador-def}. Réciproquement, chaque sous-ensemble de~\eqref{ix:eq:ch-operador-def} est donné par un code de formule et un tuple fini de paramètres ; d'après l'hypothèse de récurrence, ces paramètres possèdent des termes et~\eqref{ix:eq:ch-terminos} forme le terme qui le désigne. À une limite, les deux constructions prennent la même réunion.
\end{proof}

\Needspace{9\baselineskip}
L’encodeur d’états est compatible avec les troncatures et envoie chaque
prolongement TPK sur un terme successeur. Soit
\(\mathcal T_N^{\rm st}\subset\mathcal T^{\rm HMT}\) la sous-classe de termes
qui code des états complets de profondeur \(N\), et soit
\(\tau_{N+1,N}:\mathcal T_{N+1}^{\rm st}\to\mathcal T_N^{\rm st}\) l’effacement
du dernier bloc avec la troncature correspondante du registre généalogique,
de la retenue et de la mémoire. Si \(J_N:X_N\to\mathcal T_N^{\rm st}\) est l’encodeur
complet, la compatibilité de \(\operatorname{Ext}_N\) prouve l’identité
correctement typée

\[
 \forall x\in X_N\ \forall y\in\operatorname{Ext}_N(x)\qquad
 \tau_{N+1,N}\bigl(J_{N+1}(y)\bigr)=J_N(x).
\tag{37e$_0$}\label{ix:eq:ch-naturalidad-codificador}
\]

Définissons au sein de la sémantique
\[
 \operatorname{Succ}^{\rm TPK}_{N+1}(J_N(x)):=
 \{t\in\mathcal T_{N+1}^{\rm st}:\exists y\in X_{N+1}\,
   (y\in\operatorname{Ext}_N(x)\ \wedge\ t=J_{N+1}(y))\}.
\]
Choisissons \(\beta\) tel que \(\mathfrak G_\beta\) contienne \(x\), \(X_N\),
\(X_{N+1}\), les graphes finis de
\(\operatorname{Ext}_N,J_N,J_{N+1}\) et leurs codes dans \(h\). Alors la
formule primitive récursive du graphe de \(\operatorname{Ext}_N\) donne

\[
 J_{N+1}[\operatorname{Ext}_N(x)]
 =\operatorname{Succ}^{\rm TPK}_{N+1}(J_N(x)),
 \qquad
 J_{N+1}[\operatorname{Ext}_N(x)]
 \in\mathfrak G_{\beta+1},
\tag{37e}\label{ix:eq:ch-sucesores-semanticos}
\]

et les deux membres sont formés par la formule de graphe de la même extension
TPK. Les équations~\eqref{ix:eq:ch-naturalidad-codificador}--\eqref{ix:eq:ch-sucesores-semanticos} sont le lien sémantique explicite : la
sous-classe des termes successeurs est exactement l’image des prolongements
admissibles et leur troncature commute. On n’identifie pas la fibre complète de
\(\rho\) à \(\operatorname{Ext}_N\), et l’on n’identifie pas
\(\operatorname{Def}\) à TPK par nomenclature.

\begin{lemma}[Transitivité, axiomes et bon ordre du modèle généalogique]
\label{ix:thm:ch-6b}
\(\mathfrak G_{\rm HMT}(h,m_\infty)\) est une classe transitive, contient tous les ordinaux, satisfait la théorie des ensembles de Zermelo--Fraenkel (ZF) et possède un bon ordre global définissable ; elle satisfait donc le choix et tous les axiomes utilisés dans la comparaison cardinale ultérieure.

\end{lemma}
\begin{proof}
Abrégeons \(B=\mathfrak G_0\) et \(u=u_{h,m}\). La récurrence
sur~\eqref{ix:eq:ch-clausura-transfinita} prouve simultanément que chaque niveau est transitif et que
\(\mathfrak G_\alpha\in\mathfrak G_{\alpha+1}\) ; si tous les ordinaux
inférieurs à \(\alpha\) sont présents, leur ensemble est définissable au niveau
correspondant et publie \(\alpha\). L’union contient donc tous les
ordinaux. L’Extensionalité et la Fondation sont héritées de l’univers ambiant ;
le Vide et l’Infini sont dans \(B\). Une fois choisi un niveau contenant les paramètres,
les formules qui définissent \(\{a,b\}\), \(\bigcup a\) et les sous-ensembles par
Séparation les publient dans un successeur.

Pour la collection--remplacement, soit \(F\) une relation fonctionnelle définissable dans \(\mathfrak G\) sur l'ensemble \(a\). Le schéma de réflexion relatif à la hiérarchie définissable \((\mathfrak G_\xi,\in,u_{h,m})_{\xi\in\mathrm{Ord}}\), appliqué dans la métathéorie ambiante à la liste finie de formules qui définit \(F\), produit un niveau \(\mathfrak G_\theta\) qui contient \(a\), les paramètres, et qui est correct pour ces formules. Les rangs des témoins de \(F[a]\) sont alors bornés par \(\theta\) ; la séparation sur \(\mathfrak G_\theta\) forme l'image dans \(\mathfrak G_{\theta+1}\). Pour l'ensemble des parties, on ne présuppose pas l'ensemble interne que l'on veut construire. Une fois \(a\in\mathfrak G_\gamma\) fixé, on considère dans la métathéorie ambiante l'ensemble déjà existant \(\mathcal P^V(a)\). Pour chaque \(b\in\mathcal P^V(a)\cap\mathfrak G\), soit \(r(b)\) sa première étape généalogique. Comme \(\mathfrak G\) est une classe définissable, la séparation dans la métathéorie ambiante forme d'abord l'ensemble \(\mathcal P^V(a)\cap\mathfrak G\). Le remplacement \textbf{dans la métathéorie ambiante} appliqué à cet ensemble borne les ordinaux \(r(b)\) par un certain \(\theta\). Par conséquent,
\[
 \mathcal P^{\mathfrak G}(a)
 =\{b\in\mathfrak G_\theta:b\subseteq a\},
\]
et le membre de droite est publié par Séparation dans
\(\mathfrak G_{\theta+1}\). Il n’introduit pas dans \(\mathfrak G\) les sous-ensembles
extérieurs à la sémantique~\eqref{ix:eq:ch-objetos-semanticos}. Tel est l’argument non circulaire de
majoration par les rangs.

Enfin, on attribue à chaque objet son premier niveau de naissance et son
plus petit \textbf{nom de dérivation} : un arbre de dérivation bien fondé et
à ramifications finies (et, par~\eqref{ix:eq:ch-terminos}, fini pour chaque nom individuel), étiqueté par un
ordinal de niveau, un code de formule et le tuple des noms de ses paramètres.
Par récurrence transfinie, on ordonne ces noms : d’abord l’étiquette ordinale,
puis le code de formule, puis, lexicographiquement, les noms des
paramètres déjà ordonnés. Le plus petit nom de chaque valeur définit une relation de
classe \(<_{\mathfrak G}\) qui compare tous les objets. C’est un bon ordre
global définissable qui, en prenant le plus petit élément de chaque ensemble non vide,
prouve le Choix.
\end{proof}


% BEGIN OWNER /Users/ruben/Documents/New project/output/REVISION_ARGUMENTAL_APERTURAS_CIERRES_20260915/IX/ES/sections/ch_elevacion_catalogal.tex
\subsection{Relèvement TPK de la fibre ternaire et partition cylindrique complète}
\label{ix:subsec:elevacion-catalogal-ch}

La fibre de \(729=3^6\) sous-cylindres de raffinement ne provient pas du quotient de transfert
\(\mathcal K_{\rm ph}\). Cet opérateur est une publication ultérieure à mémoire
réduite et ne peut pas engendrer d’histoires. Le recouvrement est construit directement
dans le système \((H_k,\rho_k,\operatorname{Ext}_k)\) au moyen des trois
prolongements d’un tour de l’holonomie nonadique.


% BEGIN OWNER /Users/ruben/Documents/New project/output/REVISION_ARGUMENTAL_APERTURAS_CIERRES_20260915/IX/ES/sections/ch_capacidad_hensel.tex
\subsubsection{Inversion de Hensel dans la tour arithmétique et réalisation par histoires}
\label{ix:sec:ch-capacidad-hensel}

Avant de relever les branches en histoires enrichies, il convient de retrouver
l’argument fini qui détermine la capacité de la tour arithmétique. On
incorpore la preuve complète du transducteur non filtré de la matrice ; son
domaine est distingué de l’arbre survivant.

\begin{proposition}[Conjugaison de la tour de Hensel non filtrée]
\label{ix:prop:ch-hensel-bruto}
Soient \(p\) premier, \(d\geq1\) et \(A\in\operatorname{GL}_d(\mathbb Z)\). Pour des vecteurs lignes, en prenant des représentants de résidus dans \(\{0,\ldots,p-1\}^d\), on définit
\[
 x_tA=u_t+px_{t+1},\qquad u_t=x_tA\bmod p.
\]
L'application
\[
 \Psi_T:(\mathbb Z/p^T\mathbb Z)^d\longrightarrow
             (\mathbb F_p^d)^T,\qquad
 x_0\longmapsto(u_0,\ldots,u_{T-1})
\]
est bijective, d’inverse
\[
 x_0\equiv\sum_{t=0}^{T-1}p^tu_tA^{-(t+1)}\pmod{p^T}.
\]
Les bijections commutent avec les réductions et donnent un homéomorphisme
\(\Psi:\mathbb Z_p^d\to(\mathbb F_p^d)^{\mathbb N}\), où
\(\mathbb Z_p:=\varprojlim_T\mathbb Z/p^T\mathbb Z\). L’actualisation
\(x_t\mapsto x_{t+1}\) correspond à supprimer le premier bloc.
\end{proposition}

\begin{proof}
L'inverse entier de \(A\) existe car son déterminant est \(\pm1\). En itérant \(x_t=(u_t+px_{t+1})A^{-1}\), on obtient
\[
 x_0=\sum_{t=0}^{T-1}p^tu_tA^{-(t+1)}+p^Tx_TA^{-T}.
\]
Modulo \(p^T\), un mot détermine une unique classe initiale. Étant donné un
mot arbitraire, fixer \(x_T=0\) et reconstruire à rebours produit une
classe qui l’émet, ce qui prouve la surjectivité. La même formule
prouve la compatibilité lors de la troncature. Le passage à la limite inverse donne
l’homéomorphisme et la relation avec le décalage du ruban.
\end{proof}

Pour \(p=3\) et \(d=6\), l’alphabet de blocs a \(3^6=729\)
éléments. Cette capacité arithmétique n’est pas identifiée par définition à
la réalisabilité de chaque bloc après chaque préfixe survivant. Le
lien avec les arêtes \(\operatorname{Ext}\) est précisément la fonction
du lemme~\ref{ix:lem:elevacion-catalogal-nonadica}. Les deux
opérations que distingue la matrice sont ainsi maintenues : représentation dans la tour brute et
réalisation sous les compatibilités APP--TRIT--TPK.

% END OWNER



% BEGIN OWNER /Users/ruben/Documents/New project/output/REVISION_ARGUMENTAL_APERTURAS_CIERRES_20260915/IX/ES/sections/ch_tres_vueltas_ext.tex
\subsubsection{Réalisation arête par arête des trois prolongements nonadiques}
\label{ix:subsubsec:tres-vueltas-ext-ch}

Cette sous-section spécifie le modèle autonome engendré par la relation
d’extension TPK qui intervient dans la preuve cardinale. On n’infère pas l’existence
d’une arête par coïncidence de types : on construit ses états initial et final,
sa relation APP, son transport de fibre et la mise à jour unaire du registre.
La construction utilise exclusivement les deux feuillets d’APP, le sélecteur TRIT
et la composition \(\mathsf{Sel}\)--\(\mathsf{Tra}\)--\(\operatorname{Upd}\) du TPK.

\paragraph{Cellules complètes et trois réductions équilibrées.}
Dans la carte \(D_9=\{1,\ldots,9\}\), nous maintiendrons toujours la cellule APP complète
\[
 a_{p,q}:=(p,q;R^\Sigma_{p,q},q^\Sigma_{p,q},
                 R^\Pi_{p,q},q^\Pi_{p,q}),
\]
avec
\[
 p+q=R^\Sigma_{p,q}+9q^\Sigma_{p,q},\qquad
 pq=R^\Pi_{p,q}+9q^\Pi_{p,q}.
\]
Le feuillet actif est enregistré dans une coordonnée séparée ; \(a_{p,q}\) n'est pas scindé en une prétendue cellule additive et une autre multiplicative. Pour \(r\in D_9\), soient
\begin{equation}
 a(r):=
 \begin{cases}
  a_{1,9},&r=1,\\
  a_{1,r-1},&2\le r\le9,
 \end{cases}
 \qquad
 \tau(r):=M_{\rm ph}(r),
 \qquad
 m(r):=\tau(r)\in\{-1,0,+1\}.
 \label{ix:eq:celula-completa-fase-ch}
\end{equation}
Le reste additif de \(a(r)\) fournit ici un représentant APP de la
phase ; cela ne transforme pas \(\Sigma\) en feuillet actif universel. Le sélecteur
dynamique active \(\Sigma\) dans \(r\in\{1,4,7\}\), active \(\Pi\) dans
\(r\in\{2,5,8\}\) et conserve les deux feuillets, sans déplacer le curseur, dans
\(r\in\{3,6,9\}\). Dans les trois cas, on transporte conjointement
\(R^\Sigma,q^\Sigma,R^\Pi,q^\Pi\). On a
\[
 (m(1),\ldots,m(9))=(1,-1,0,1,-1,0,1,-1,0).
\]
Les trois paires ordonnées par la base nonadique sont
\begin{equation}
 \mathfrak p_{+1}:=(1,4),\qquad
 \mathfrak p_{-1}:=(2,5),\qquad
 \mathfrak p_{0}:=(3,6).
 \label{ix:eq:tres-pares-internos-ch}
\end{equation}
Ces paires sont la section normalisée qui prend les deux premières occurrences de chaque type TRIT sous le déplacement \(r\mapsto r+3\) ; le troisième bloc de phases \(7,8,9\) conserve le contrôle de clôture du tour. Leurs deux membres appartiennent à une même fibre de \(M_{\rm ph}\). La réduction équilibrée du chapitre TRIT donne, sans introduire de coefficient cible,
\begin{equation}
 \begin{array}{c|c|c|c}
 \varepsilon&m(\mathfrak p_\varepsilon^{(1)})
                  +m(\mathfrak p_\varepsilon^{(2)})
   &r_3\bigl(2\varepsilon\bigr)&
   \chi(\varepsilon,\varepsilon)\\ \hline
 +1&1+1=(-1)+3\cdot1&-1&+1\\
 0&0+0=0+3\cdot0&0&0\\
 -1&(-1)+(-1)=1+3\cdot(-1)&+1&-1
 \end{array}
 \label{ix:eq:tres-carry-derivados-ch}
\end{equation}
Par conséquent,
\begin{equation}
 \chi(\varepsilon,\varepsilon)
 =\frac{2\varepsilon-r_3(2\varepsilon)}3=\varepsilon.
 \label{ix:eq:carry-no-parametrico-ch}
\end{equation}
Le signe de la branche est ainsi la sortie du cocycle équilibré appliqué à deux
émissions TRIT produites par APP. Chaque paire fournit l’une des trois
sorties de retenue. Par souci de clarté, nous indexerons ensuite la paire par
\(\varepsilon(\mathfrak p):=
\chi(m(\mathfrak p^{(1)}),m(\mathfrak p^{(2)}))\) ; le symbole
\(\varepsilon\) ne fournit pas à \(\operatorname{Upd}^{\rm cal}\) une retenue libre.

L’architecture de phase, le relèvement équilibré et le transport de
mémoire sont réalisés ici au moyen de vingt-sept incidences locales et de leurs
graphes d’extension. La section normalisée fixe des représentants des
fibres de \(M_{\rm ph}\) ; la couverture qui sera démontrée utilise l’existence
de ces représentants et demeure invariante lorsqu’on les remplace par une autre
section compatible.

\paragraph{Données typées des neuf arêtes.}
Soit
\(K_{\rm Carr}:=\langle\kappa_{\rm Carr}\rangle_{\rm Ab}\) le module de
retenue. Les actions sont typées par
\[
 H_{\rm Carr}:\mathscr P_{\rm TPK}^{\rm op}
   \longrightarrow\operatorname{End}(K_{\rm Carr}),\qquad
 Y_{\rm Carr}:\mathscr P_{\rm TPK}
   \longrightarrow\operatorname{End}(K_{\rm Carr}).
\]
Dans ce modèle engendré, on prend les actions triviales \(H_{\rm Carr}(e)=Y_{\rm Carr}(e)=\operatorname{id}_{K_{\rm Carr}}\) : le feuillet actif et le chemin sont conservés comme coordonnées propres du registre et n'agissent pas sur le terme libre de retenue. Soit en outre \((\mathsf M_d,\jmath_{d',d})\) le système inductif gradué de mémoire. Pour ne pas confondre deux représentants identifiés par la colimite, l'état conserve en outre le degré comme coordonnée et l'on utilise l'espace porteur étiqueté
\[
 \widetilde{\mathsf M}:=
 \bigsqcup_{d\ge1}\bigl(\{d\}\times\mathsf M_d\bigr),
 \qquad
 \operatorname{gr}_{\mathsf M}(d,\mu)=d.
\]
La publication dans la mémoire limite est donnée par
\[
\begin{aligned}
 \mathsf M&:=\varinjlim_d\mathsf M_d,&
 \jmath_d&:\mathsf M_d\longrightarrow\mathsf M,\\
 u&:=\jmath_d(u_d),&
 \iota_{\mathsf M}:\mathbb Z&\longrightarrow\mathsf M,
 \qquad n\longmapsto nu .
\end{aligned}
\]
Ici, \(\jmath_d\) est l’application canonique de la colimite et les \(u_d\) forment une
famille compatible. Sa composante libre permet
le lecteur \(\deg_u(nu)=n\).
Les projections de configuration, de phase, de mémoire, de retenue, de route et de
survie seront typées sur le porteur calibré construit
récursivement ci-dessous ; aucun espace ambiant ne sera importé pour elles.
Pour chaque arête construite \(e\), la
composante calibrée
\(\operatorname{Upd}^{\rm cal}_e:=\mathsf{Mem}^{\rm cal}_e\) désignera la
troisième application de la factorisation ; sa loi est fixée champ par champ plus
bas et n’est pas présupposée comme restriction d’une relation ambiante.

Fixons la profondeur initiale \(k_0=1\), une configuration complète
\(q_\star\in\mathcal Q_{\rm obs}\) de phase un et
\[
 q_{n,r}:=F_{\rm obs}^{9n+r-1}(q_\star),\qquad
 q_{n,10}:=q_{n+1,1}.
\]
Comme le calendrier observable est périodique, sa valeur \(q_{n,r}\) ne
distingue pas à elle seule des apparitions séparées par douze tours. Le modèle
calibré emploie donc le relèvement étiqueté
\[
 \widetilde q_{n,r}:=(q_{n,r},n,r),\qquad
 \widetilde q_{n,10}:=\widetilde q_{n+1,1},\qquad
 \pi_{\rm obs}(\widetilde q_{n,r})=q_{n,r}.
\]
La dynamique élevée \(\widetilde F_{\rm obs}(\widetilde q_{n,r})
=\widetilde q_{n,r+1}\) se projette sur \(F_{\rm obs}\) ; l'étiquette ne modifie pas la phase physique, mais distingue ses niveaux d'apparition. On étiquette la copie de la cellule à chaque niveau au moyen de
\[
 a_{n,r}:=(n,r;a(r)),\qquad
 (n,r)^+:=
 \begin{cases}
  (n,r+1),&1\le r<9,\\
  (n+1,1),&r=9.
 \end{cases}
\]
L’étiquette évite seulement d’identifier deux apparitions d’une même cellule ; les
deux divisions euclidiennes et leurs quatre coordonnées sont celles de \(a(r)\).
Pour le tour \(n\ge0\) et la phase \(r\in D_9\), le symbole
\[
 e_{\varepsilon,n,r}:\widetilde q_{n,r}
 \longrightarrow\widetilde q_{n,r+1},
 \qquad r^+:=1+(r\bmod9),
\]
occupe le niveau \(d_{n,r}:=k_n+r-1\), où \(k_n:=k_0+9n\). Son registre local contient la cellule complète \(a(r)\), le feuillet actif déterminé par \(M_{\rm ph}(r)\), le type \(\tau(r)\), la marque d'appartenance à \(\mathfrak p_\varepsilon\) et, lorsqu'il atteint le second membre de cette paire, l'identité certifiée \eqref{ix:eq:tres-carry-derivados-ch}. Nous définissons ses opérateurs uniquement à partir de ce registre :
\begin{align}
 \mathsf C_{\mathsf M}(e_{\varepsilon,n,r})
   &=\jmath_{d_{n,r}+1,d_{n,r}},&
 \kappa_{\mathsf M}(e_{\varepsilon,n,r})
   &=\mathbf1_{\{r=9\}}u_{d_{n,r}+1},
 \label{ix:eq:memoria-arista-ch}\\
       \operatorname{Carr}(e_{\varepsilon,n,r})
   &=d_{\varepsilon,r}\kappa_{\rm Carr},&
 d_{\varepsilon,r}
   &:=\mathbf1_{\{r=\mathfrak p_\varepsilon^{(2)}\}}
       \chi\!\left(m(\mathfrak p_\varepsilon^{(1)}),
                    m(\mathfrak p_\varepsilon^{(2)})\right).
 \label{ix:eq:carry-arista-derivado-ch}
\end{align}
Pour les identités et pour toute composition admissible, on définit en outre
\begin{align}
 \mathsf C_{\mathsf M}(\operatorname{id}_{\widetilde q};d)
   &:=\operatorname{id}_{\mathsf M_d},&
 \kappa_{\mathsf M}(\operatorname{id}_{\widetilde q};d)&:=0,
 \label{ix:eq:memoria-identidad-ch}\\
 \mathsf C_{\mathsf M}(e_t\cdots e_1)
   &:=\mathsf C_{\mathsf M}(e_t)\circ\cdots\circ
      \mathsf C_{\mathsf M}(e_1),&
 \kappa_{\mathsf M}(e_2e_1)
   &:=\mathsf C_{\mathsf M}(e_2)
        \bigl(\kappa_{\mathsf M}(e_1)\bigr)
      +\kappa_{\mathsf M}(e_2).
 \label{ix:eq:cociclo-memoria-calibrado-ch}
\end{align}
L’identité porte le degré \(d\) car une même configuration observable
peut réapparaître à des profondeurs distinctes. La dernière égalité, itérée, définit
\(\kappa_{\mathsf M}(e_t\cdots e_1)\) pour toute longueur et est la loi de
cocycle du système direct ; les inclusions satisfont
\(\jmath_{d+2,d+1}\jmath_{d+1,d}=\jmath_{d+2,d}\).
On fixe en outre
\(s(e_{\varepsilon,n,r})=\top\). Les vingt-sept entrées suivantes
certifient les coordonnées qui distinguent les trois prolongements ; les champs
universels du transport sont spécifiés immédiatement après. Les symboles \(I\) et \(II\)
désignent respectivement la première émission conservée et la seconde émission
qui effectue la réduction ; un tiret indique un accroissement de retenue nul, mais non
une absence d’arête.
\begin{center}
\small
\setlength{\tabcolsep}{3.2pt}
\begin{tabular}{c c cc cc cc c}
\toprule
\(r\)&\(M_{\rm ph}(r)\)&\multicolumn{2}{c}{\(\mathfrak p_{+1}\)}&
\multicolumn{2}{c}{\(\mathfrak p_0\)}&
\multicolumn{2}{c}{\(\mathfrak p_{-1}\)}&
\(\Delta\operatorname{mem}\)\\
& &marque&\(d_{+1,r}\)&marque&\(d_{0,r}\)&marque&\(d_{-1,r}\)&\\
\midrule
1&\(+1\)&I&0&--&0&--&0&0\\
2&\(-1\)&--&0&--&0&I&0&0\\
3&0&--&0&I&0&--&0&0\\
4&\(+1\)&II&\(+1\)&--&0&--&0&0\\
5&\(-1\)&--&0&--&0&II&\(-1\)&0\\
6&0&--&0&II&0&--&0&0\\
7&\(+1\)&--&0&--&0&--&0&0\\
8&\(-1\)&--&0&--&0&--&0&0\\
9&0&--&0&--&0&--&0&\(u\)\\
\bottomrule
\end{tabular}
\end{center}
Ainsi, \(d_{\varepsilon,r}\) est calculé à partir de deux sorties du sélecteur et
est distinct des quotients \(q^\Sigma,q^\Pi\) de la cellule. Le registre
conserve également la paire marquée et le résidu équilibré ; en particulier, la
branche nulle n’est pas confondue avec une absence d’événement.

La loi de composition de la retenue employée par la catégorie TPK est
\begin{equation}
 \operatorname{Carr}(\operatorname{id}_{\widetilde q})=0,\qquad
 \operatorname{Carr}(e_2e_1)
 =H_{\rm Carr}(e_1)\bigl(\operatorname{Carr}(e_2)\bigr)
  +Y_{\rm Carr}(e_2)\bigl(\operatorname{Carr}(e_1)\bigr).
 \label{ix:eq:ley-carry-correcta-ch}
\end{equation}
Sur les identités et les chemins, ils se prolongent fonctoriellement :
\[
\begin{aligned}
 H_{\rm Carr}(\operatorname{id})
   &=Y_{\rm Carr}(\operatorname{id})=\operatorname{id},\\
 H_{\rm Carr}(e_t\cdots e_1)
   &=H_{\rm Carr}(e_1)\circ\cdots\circ H_{\rm Carr}(e_t),\\
 Y_{\rm Carr}(e_t\cdots e_1)
   &=Y_{\rm Carr}(e_t)\circ\cdots\circ Y_{\rm Carr}(e_1).
\end{aligned}
\]
Dans ce modèle, toutes ces actions restent l'identité, mais leur typage explicite quel endomorphisme agit sur chaque terme de \eqref{ix:eq:ley-carry-correcta-ch}.
Sur les arêtes calibrées, \(H_{\rm Carr}=Y_{\rm Carr}=\operatorname{id}\), si bien que la composition additionne les accroissements déjà produits. L'équation \eqref{ix:eq:ley-carry-correcta-ch}, et non une identification entre orientation et quotient, gouverne le registre accumulé.

\paragraph{Construction indépendante des tuples de transition.}
Un état de la fibre sera écrit, dans l’ordre de ses champs pertinents,
\[
 x=(a;q;\tau,o,h;b_{\rm ref};
       R^\Sigma,q^\Sigma,R^\Pi,q^\Pi;c;
       \operatorname{Sig};C,\partial C;p_{\APP},p_{\TRIT};
       \lambda;\operatorname{gr}_{\mathsf M};\operatorname{mem};
       s;\operatorname{Led}).
\]
La notation de mise à jour \(x[f\leftarrow v]\) ne change que le champ
indiqué. Pour \(h\in\{(\Sigma,\Pi),(\Pi,\Sigma)\}\), définissons
\[
 h_r(h)=
 \begin{cases}
  (\Sigma,\Pi),&r\in\{1,4,7\},\\
  (\Pi,\Sigma),&r\in\{2,5,8\},\\
 h,&r\in\{3,6,9\}.
 \end{cases}
\]
La coordonnée complète de feuillet et d’arithmétique est abrégée par
\[
 \widehat h(x):=
 \bigl(h(x),R^\Sigma(x),q^\Sigma(x),R^\Pi(x),q^\Pi(x)\bigr).
\]
Soit \(\widehat{\mathcal H}^{\rm cal}_{\widetilde q_{n,r}}\) l’ensemble des
quintuplets
\[
 \bigl(h,R^\Sigma_{a_{n,r}},q^\Sigma_{a_{n,r}},
          R^\Pi_{a_{n,r}},q^\Pi_{a_{n,r}}\bigr),
 \qquad
 h\in\{(\Sigma,\Pi),(\Pi,\Sigma)\},
\]
dont les quatre données arithmétiques sont les deux divisions euclidiennes de la cellule
APP étiquetée \(a_{n,r}\). Cette définition directe ne présuppose pas d’état
sur un porteur qui n’est pas encore construit. Soient
\(\mathsf{Hist}^{\rm cal}_{\widetilde q}\) l’ensemble des chemins finis
admissibles du graphe d’arêtes étiquetées qui aboutissent à
\(\widetilde q\), y compris l’identité, et
\(\mathsf B^{\rm cal}_{\widetilde q}\) l’ensemble de leurs suites
ordonnées d’extrémités. Nous définissons
\[
 \mathsf S^{\rm cal}_{\widetilde q}:=
 \widehat{\mathcal H}^{\rm cal}_{\widetilde q}
 \times\{o_{\rightarrow},o_{\circ},o_{\leftarrow}\}
 \times K_{\rm Carr}\times
 \mathsf{Hist}^{\rm cal}_{\widetilde q}\times
 \mathsf B^{\rm cal}_{\widetilde q}\times\mathsf M .
\]
La signature calibrée n’est pas postulée : c’est l’application de réunion typée
\[
 \operatorname{Sig}^{\rm cal}_{\widetilde q}:
 \widehat{\mathcal H}^{\rm cal}_{\widetilde q}
 \times\{o_{\rightarrow},o_{\circ},o_{\leftarrow}\}
 \times K_{\rm Carr}\times
 \mathsf{Hist}^{\rm cal}_{\widetilde q}\times
 \mathsf B^{\rm cal}_{\widetilde q}\times\mathsf M
 \longrightarrow\mathsf S^{\rm cal}_{\widetilde q},
\]
\[
 (\widehat h,o,c,\lambda,\partial C,\mu)
 \longmapsto(\widehat h,o,c,\lambda,\partial C,\mu).
\]
De même, \(\pi_{\mathcal F}\) désigne la projection qui élimine uniquement les deux premières coordonnées \((a;q)\) du tuple d'état ; elle conserve \(\tau,o,h,b_{\rm ref}\), les quatre données euclidiennes, la retenue, la signature, le cylindre, la frontière, les préfixes, le chemin, le degré de mémoire, la mémoire, la survie et le registre.
L'orientation correspondante est fixée par \(o(+1)=o_{\rightarrow}\), \(o(0)=o_{\circ}\) et \(o(-1)=o_{\leftarrow}\). La fibre APP calibrée sur \(\widetilde q_{n,r}\) contiendra la copie étiquetée \(a_{n,r}\) déjà définie. L'état initial est fixé sans coordonnées libres :
\[
 \widehat h_\star:=
 \bigl((\Sigma,\Pi),R^\Sigma_{1,9},q^\Sigma_{1,9},
       R^\Pi_{1,9},q^\Pi_{1,9}\bigr).
\]
\begin{equation}
\begin{aligned}
 x_\star:=\bigl(&a_{0,1};\widetilde q_{0,1};
 +1,o_{\rightarrow},(\Sigma,\Pi);\bot;
 R^\Sigma_{1,9},q^\Sigma_{1,9},R^\Pi_{1,9},q^\Pi_{1,9};0;\\
 &\operatorname{Sig}^{\rm cal}_{\widetilde q_{0,1}}
       (\widehat h_\star,o_{\rightarrow},0,
        \operatorname{id}_{\widetilde q_{0,1}},\varnothing,
        \jmath_{k_0}(0_{\mathsf M_{k_0}}));
 \varnothing,\varnothing;
 (a_{0,1}),(+1);\operatorname{id}_{\widetilde q_{0,1}};
 k_0;0_{\mathsf M_{k_0}};\top;\varnothing\bigr).
\end{aligned}
\label{ix:eq:estado-inicial-calibrado-ch}
\end{equation}
Le tuple ne contient pas de coordonnées indéterminées et satisfait
\(s(x_\star)=\top\). Soit \(G_0:=\{x_\star\}\).
Supposant \(G_n\) construit, fixons \(x\in G_n\),
\(\varepsilon\in\{-1,0,+1\}\) et
\(x_{\varepsilon,0}(x):=x\). Pour \(r=1,\ldots,9\), on construit d’abord
le registre élémentaire
\begin{equation}
\begin{aligned}
 \ell_{\varepsilon,n,r}:=\bigl(&e_{\varepsilon,n,r},a_{n,r},a_{(n,r)^+},
 \widetilde q_{n,r},\widetilde q_{n,r+1},M_{\rm ph}(r),\\
 &h_r(h(x_{\varepsilon,r-1}(x))),o(M_{\rm ph}(r)),\mathfrak p_\varepsilon,
 d_{\varepsilon,r},\\
 &\mathsf C_{\mathsf M}(e_{\varepsilon,n,r}),
 \kappa_{\mathsf M}(e_{\varepsilon,n,r}),
 \operatorname{Carr}(e_{\varepsilon,n,r})\bigr),
\end{aligned}
 \label{ix:eq:registro-elemental-ext-ch}
\end{equation}
et on pose
\(\operatorname{Led}^{\rm new}_{\varepsilon,n,r}:=
\operatorname{Led}(x_{\varepsilon,r-1}(x))^\frown
\ell_{\varepsilon,n,r}\). On construit alors le tuple
\begin{equation}
 y_{\varepsilon,r}(x):=
 x_{\varepsilon,r-1}(x)
 [\tau\leftarrow M_{\rm ph}(r),\
  h\leftarrow h_r(h(x_{\varepsilon,r-1}(x))),\
  o\leftarrow o(M_{\rm ph}(r)),\
  b_{\rm ref}\leftarrow\mathfrak p_\varepsilon],
 \label{ix:eq:tupla-sel-cal-ch}
\end{equation}
où la sélection laisse intacts la cellule, les quotients, la retenue, la route, la mémoire,
le cylindre, la frontière, la survie et le registre ; la coordonnée
\(b_{\rm ref}\) explicite la réduction équilibrée choisie. On ne présuppose pas ici
l’appartenance à une relation ambiante.

On définit ensuite \(z_{\varepsilon,r}(x)\) champ par champ : sa
configuration est \(\widetilde q_{n,r+1}\) ; sa cellule complète est
\(a_{(n,r)^+}\), avec les
quatre résidus et quotients recalculés par les deux divisions euclidiennes ;
il conserve le feuillet sélectionné et l’orientation ; et il effectue les six
mises à jour suivantes, dans lesquelles seul l’argument commun \(x\) est omis :
\begin{equation}
\begin{aligned}
 C(z_{\varepsilon,r})&:=C(y_{\varepsilon,r})^\frown
       e_{\varepsilon,n,r},\\
 \partial C(z_{\varepsilon,r})&:=\partial C(y_{\varepsilon,r})^\frown
       (\widetilde q_{n,r},\widetilde q_{n,r+1}),\\
 p_{\APP}(z_{\varepsilon,r})&:=p_{\APP}(y_{\varepsilon,r})^\frown
       a_{(n,r)^+},\\
 p_{\TRIT}(z_{\varepsilon,r})&:=p_{\TRIT}(y_{\varepsilon,r})^\frown
       \tau(r),\\
 \operatorname{Led}(z_{\varepsilon,r})&:=
       \operatorname{Led}^{\rm new}_{\varepsilon,n,r},\\
 \lambda(z_{\varepsilon,r})&:=e_{\varepsilon,n,r}
       \lambda(y_{\varepsilon,r}),
\end{aligned}
 \label{ix:eq:tupla-tra-cal-ch}
\end{equation}
et maintient encore les champs \(c,\operatorname{gr}_{\mathsf M},\operatorname{mem},s\) de \(y_{\varepsilon,r}(x)\). Sa signature, en revanche, est à nouveau typée dans la configuration de destination au moyen de la présignature transportée
\begin{equation}
\begin{aligned}
 \operatorname{Sig}(z_{\varepsilon,r}(x))
  :=\operatorname{Sig}^{\rm cal}_{\widetilde q_{n,r+1}}
       \Bigl(&\widehat h(z_{\varepsilon,r}(x)),
             o(z_{\varepsilon,r}(x)),c(y_{\varepsilon,r}(x)),
             \lambda(z_{\varepsilon,r}(x)),\\[-2pt]
             &\partial C(z_{\varepsilon,r}(x)),
             \jmath_{d_{n,r}}
                (\operatorname{mem}_{d_{n,r}}
                   (y_{\varepsilon,r}(x)))\Bigr).
\end{aligned}
\label{ix:eq:prefirma-transporte-cal-ch}
\end{equation}
Ainsi, et seulement dans
\(\mathsf{Tra}^{\rm cal}\), on ajoute l’arête, ses extrémités, la marque
\((\mathfrak p_\varepsilon,r,d_{\varepsilon,r})\) et ses cocycles au registre des
incidences.

L'état initial satisfait
\[
 c(x_\star)=0=\operatorname{Carr}
   (\operatorname{id}_{\widetilde q_{0,1}})
   =\operatorname{Carr}(\lambda(x_\star)).
\]
Dans les états complets \(x_{\varepsilon,r-1}(x)\), et aussi après la
sélection \(y_{\varepsilon,r}(x)\), on conserve l’invariant typé
\(c(v)=\operatorname{Carr}(\lambda(v))\), puisque
\(\mathsf{Sel}^{\rm cal}\) ne modifie aucun de ces deux champs. Le
transport fait avancer la route mais laisse la retenue en attente ; le tuple de
préactualisation satisfait donc exactement
\[
 c(z_{\varepsilon,r}(x))=c(y_{\varepsilon,r}(x))
 =\operatorname{Carr}(\lambda(y_{\varepsilon,r}(x))).
\]
L'application \(\operatorname{Upd}^{\rm cal}\) restaure alors \(c(x_{\varepsilon,r}(x))=
\operatorname{Carr}(\lambda(x_{\varepsilon,r}(x)))\) au moyen de la loi de composition, sans attribuer cet invariant à l'état intermédiaire \(z\).

Enfin, on applique la fonction unaire
\(\operatorname{Upd}^{\rm cal}_{e_{\varepsilon,n,r}}
=\mathsf{Mem}^{\rm cal}_{e_{\varepsilon,n,r}}\) à ce tuple. Son
résultat \(x_{\varepsilon,r}(x)\) est donné par
\begin{equation}
\begin{aligned}
 c\bigl(x_{\varepsilon,r}(x)\bigr)
   &=\operatorname{Carr}\bigl(\lambda(z_{\varepsilon,r}(x))\bigr)\\
   &=H_{\rm Carr}(\lambda(y_{\varepsilon,r}(x)))
       \bigl(\operatorname{Carr}(e_{\varepsilon,n,r})\bigr)
     +Y_{\rm Carr}(e_{\varepsilon,n,r})
       \bigl(c(y_{\varepsilon,r}(x))\bigr),\\
 \lambda\bigl(x_{\varepsilon,r}(x)\bigr)
   &=\lambda(z_{\varepsilon,r}(x)),\\
 \operatorname{gr}_{\mathsf M}
      \bigl(x_{\varepsilon,r}(x)\bigr)
   &=d_{n,r}+1,\\
 \operatorname{mem}_{d_{n,r}+1}
       \bigl(x_{\varepsilon,r}(x)\bigr)
   &=\mathsf C_{\mathsf M}(e_{\varepsilon,n,r})
       \operatorname{mem}_{d_{n,r}}(z_{\varepsilon,r}(x))
     +\kappa_{\mathsf M}(e_{\varepsilon,n,r}),\\
 s\bigl(x_{\varepsilon,r}(x)\bigr)
   &=s(e_{\varepsilon,n,r})\wedge s(z_{\varepsilon,r}(x)),\\
 \operatorname{Sig}\bigl(x_{\varepsilon,r}(x)\bigr)
   &=\operatorname{Sig}^{\rm cal}_{\widetilde q_{n,r+1}}
       \Bigl(\widehat h(z_{\varepsilon,r}(x)),
             o(z_{\varepsilon,r}(x)),
             c(x_{\varepsilon,r}(x)),
             \lambda(z_{\varepsilon,r}(x)),
             \partial C(z_{\varepsilon,r}(x)),\\[-2pt]
   &\hspace{112pt}
             \jmath_{d_{n,r}+1}
              \bigl(\operatorname{mem}_{d_{n,r}+1}
                 (x_{\varepsilon,r}(x))\bigr)\Bigr).
\end{aligned}
\label{ix:eq:upd-unaria-tres-vueltas-ch}
\end{equation}
Tous les champs non énumérés dans \eqref{ix:eq:upd-unaria-tres-vueltas-ch} coïncident avec ceux de \(z_{\varepsilon,r}(x)\). En particulier, la signature est calculée avec le feuillet, l'orientation, la frontière et le chemin déjà fixés par \(\mathsf{Sel}^{\rm cal}\) et \(\mathsf{Tra}^{\rm cal}\), et avec la retenue et la mémoire calculées dans les lignes précédentes ; elle n'est pas définie circulairement à partir de l'état final. Ni \(\varepsilon\), ni \(d_{\varepsilon,r}\), ni \(u\) ne sont des arguments supplémentaires de \(\operatorname{Upd}^{\rm cal}\) : la fonction lit les trois données de l'arête déjà construite.

\paragraph{Récurrence simultanée des niveaux.}
Une fois les neuf étapes précédentes achevées pour un \(G_n\) déjà construit,
on définit immédiatement
\begin{equation}
 \gamma_{\varepsilon,k_n}(x):=x_{\varepsilon,9}(x),\qquad
 G_{n+1}:=\{\gamma_{\varepsilon,k_n}(x):
       x\in G_n,\ \varepsilon\in\TRIT\}.
 \label{ix:eq:injerto-total-tres-vueltas-ch}
\end{equation}
\label{ix:eq:tres-vueltas-tpk-ch}
La base \(G_0=\{x_\star\}\), les tuples intermédiaires et
\eqref{ix:eq:injerto-total-tres-vueltas-ch} constituent une unique définition par
récurrence sur \(n\). Par conséquent, aucun porteur ultérieur n’est utilisé pour
supposer l’existence de \(G_n\) : le niveau suivant n’est introduit
qu’après avoir construit toutes ses arêtes à partir du niveau précédent.

\paragraph{Relation d’extension TPK et appartenance effective.}
Soit \(\mathcal E^{\rm cal}_{\widetilde q}\) la plus petite famille étiquetée qui contient
\(x_\star\) et, à chaque étape de la récurrence précédente, les quatre tuples
\(x_{\varepsilon,r-1}(x),y_{\varepsilon,r}(x),
z_{\varepsilon,r}(x),x_{\varepsilon,r}(x)\) dans leurs configurations d’origine
et de destination. Soit également
\(\mathcal A^{\rm cal}_{\widetilde q_{n,r}}:=\{a_{n,r}\}\), et posons
\[
 \widetilde{\mathcal Q}_{\rm obs}:=
   \{\widetilde q_{n,r}:n\ge0,\ 1\le r\le9\},
 \qquad
 \mathcal E^{\rm cal}:=\bigsqcup_{\widetilde q}
     \mathcal E^{\rm cal}_{\widetilde q},
 \qquad
 \mathcal F^{\rm cal}_{\widetilde q}:=
     \pi_{\mathcal F}(\mathcal E^{\rm cal}_{\widetilde q}).
\]
Pour chaque profondeur \(d\), soit
\(\mathcal E^{\rm cal}_d:=
\operatorname{gr}_{\mathsf M}^{-1}(d)\subset\mathcal E^{\rm cal}\). La
coordonnée explicite de degré, et non la seule appartenance d’un représentant à
une image du système direct, rend ces porteurs disjoints. Sur eux,
les projections sont typées
\[
\begin{aligned}
 \operatorname{cfg}&:\mathcal E^{\rm cal}\longrightarrow
     \widetilde{\mathcal Q}_{\rm obs},\\
 g(x)&:=\operatorname{fase}\!
       \bigl(\pi_{\rm obs}(\operatorname{cfg}(x))\bigr)\in C_9,\\
 c&:\mathcal E^{\rm cal}\longrightarrow K_{\rm Carr},\\
 s&:\mathcal E^{\rm cal}\longrightarrow\{\bot,\top\},\\
 \lambda&:\mathcal E^{\rm cal}_{\widetilde q}
       \longrightarrow\mathsf{Hist}^{\rm cal}_{\widetilde q},\\
 \operatorname{gr}_{\mathsf M}&:\mathcal E^{\rm cal}
       \longrightarrow\mathbb N_{\ge1},\\
 \operatorname{mem}_d&:\mathcal E^{\rm cal}_d
       \longrightarrow\mathsf M_d,\\
 \operatorname{mem}(x)&:=
       \jmath_d\bigl(\operatorname{mem}_d(x)\bigr)\in\mathsf M
       \quad(x\text{ de profondeur }d).
\end{aligned}
\]
Ici, \(\operatorname{cfg}(x)\) est la seconde coordonnée du tuple et
\(\operatorname{fase}\) lit la phase du calendrier observable. Pour les trois
classes de porteur, on fixe les diagonales
\[
\begin{aligned}
 \Delta_{\APP,\widetilde q}
   &:=\{(a,a):a\in\mathcal A^{\rm cal}_{\widetilde q}\},\\
 \Delta_{{\rm fib},\widetilde q}
   &:=\{(f,f):f\in\mathcal F^{\rm cal}_{\widetilde q}\},\\
 \Delta_{{\rm enr},\widetilde q}
   &:=\{(x,x):x\in\mathcal E^{\rm cal}_{\widetilde q}\}.
\end{aligned}
\]
En particulier,
\[
 a_{0,1}\in\mathcal A^{\rm cal}_{\widetilde q_{0,1}},
 \qquad
 x_\star\in\mathcal E^{\rm cal}_{\widetilde q_{0,1}}.
\]
Sur ces espaces porteurs, on définit par leurs graphes complets, et non au moyen d'implications nécessaires, les relations
\begin{align}
 \mathsf{Sel}^{\rm cal}_{e_{\varepsilon,n,r}}
   &:=\{(x_{\varepsilon,r-1}(x),y_{\varepsilon,r}(x)):
          x\in G_n\},\\
 \mathsf{Tra}^{\rm cal}_{e_{\varepsilon,n,r}}
   &:=\{(y_{\varepsilon,r}(x),z_{\varepsilon,r}(x)):
          x\in G_n\},\\
 \operatorname{Ext}^{\APP,\rm cal}_{e_{\varepsilon,n,r}}
   &:=\{(a_{n,r},a_{(n,r)^+})\},
 \label{ix:eq:ext-app-calibrada-ch}\\
 \operatorname{Ext}^{\rm fib,cal}_{e_{\varepsilon,n,r}}
   &:=\{(\pi_{\mathcal F}x_{\varepsilon,r-1}(x),
          \pi_{\mathcal F}x_{\varepsilon,r}(x)):x\in G_n\}.
 \label{ix:eq:ext-fibra-calibrada-ch}
\end{align}
L'extension enrichie synchronisée est
\begin{equation}
 \operatorname{Ext}^{\rm enr,cal}_{e_{\varepsilon,n,r}}
 :=\{(x_{\varepsilon,r-1}(x),x_{\varepsilon,r}(x)):
       x\in G_n\}.
 \label{ix:eq:ext-enriquecida-calibrada-ch}
\end{equation}
La troisième application est définie par
\begin{equation}
 \operatorname{graph}
 \bigl(\operatorname{Upd}^{\rm cal}_{e_{\varepsilon,n,r}}\bigr)
 :=\{(z_{\varepsilon,r}(x),x_{\varepsilon,r}(x)):
       x\in G_n\},
 \label{ix:eq:grafo-upd-calibrado-ch}
\end{equation}
où la seconde composante est exactement celle de
\eqref{ix:eq:upd-unaria-tres-vueltas-ch}. Par conséquent,
\begin{equation}
 \mathcal U^{(1),\rm cal}_{e_{\varepsilon,n,r}}
 :=\operatorname{Upd}^{\rm cal}_{e_{\varepsilon,n,r}}\circ
   \mathsf{Tra}^{\rm cal}_{e_{\varepsilon,n,r}}\circ
   \mathsf{Sel}^{\rm cal}_{e_{\varepsilon,n,r}}
 =\{(x_{\varepsilon,r-1}(x),x_{\varepsilon,r}(x)):
      x\in G_n\}.
 \label{ix:eq:transicion-cal-en-tpk-ch}
\end{equation}
Chaque opérateur est indexé par l’arête
\(e_{\varepsilon,n,r}\) produite par la réduction
\(\mathfrak p_\varepsilon\), et non par un trit fourni de l’extérieur. La
coordonnée \(b_{\rm ref}\) distingue en outre leurs codomaines. Par conséquent,
\(\mathsf{Sel}^{\rm cal}_{e}\),
\(\mathsf{Tra}^{\rm cal}\) et \(\operatorname{Upd}^{\rm cal}\) sont des fonctions
sur leurs domaines respectifs, tandis que
l’union
\(\bigcup_{\varepsilon\in\TRIT}
\mathcal U^{(1),\rm cal}_{e_{\varepsilon,n,r}}\) conserve délibérément
les trois sorties de la trisection.

Pour \(\bullet\in\{\APP,{\rm fib},{\rm enr}\}\), les identités et les compositions ne restent pas implicites. On définit
\begin{align}
 \operatorname{Ext}^{\bullet,\rm cal}_{\operatorname{id}_{\widetilde q}}
   &:=\Delta_{\bullet,\widetilde q},\\
 \operatorname{Ext}^{\bullet,\rm cal}_{e_t\cdots e_1}
   &:=\operatorname{Ext}^{\bullet,\rm cal}_{e_t}\circ\cdots\circ
      \operatorname{Ext}^{\bullet,\rm cal}_{e_1}
 \label{ix:eq:ext-palabra-calibrada-ch}
\end{align}
pour tout mot admissible d’arêtes. Les deux clôtures de chemins APP et de
fibre sont formées séparément,
\begin{equation}
 \operatorname{Path}(\operatorname{Ext}^{\APP,\rm cal}),\qquad
 \operatorname{Path}(\operatorname{Ext}^{\rm fib,cal}),
 \label{ix:eq:dos-clausuras-ext-ch}
\end{equation}
et \(\mathsf{Tra}^{\rm cal}\) les synchronise arête par arête.

Enfin, les niveaux d'états complets et leurs troncatures sont fixés par la récurrence elle-même :
\begin{align}
 H^{\rm cal}_{d_{n,r}}
   &:=\{x_{\varepsilon,r-1}(x):
          x\in G_n,\ \varepsilon\in\TRIT\},\\
 H^{\rm cal}_{d_{n,r}+1}
   &:=\{x_{\varepsilon,r}(x):
          x\in G_n,\ \varepsilon\in\TRIT\},\\
 \rho^{\rm cal}_{d_{n,r}+1,d_{n,r}}
   \bigl(x_{\varepsilon,r}(x)\bigr)
   &:=x_{\varepsilon,r-1}(x).
 \label{ix:eq:truncamiento-calibrado-definido-ch}
\end{align}
Le dernier registre conserve \((n,r,\mathfrak p_\varepsilon)\), de sorte que
le prédécesseur du membre de gauche est unique et \(\rho^{\rm cal}\) est bien
définie. Pour \(b>a\), on pose
\(\rho^{\rm cal}_{b,a}:=\rho^{\rm cal}_{a+1,a}\circ\cdots\circ
\rho^{\rm cal}_{b,b-1}\). Cette application retrouve l’état précédent
complet, y compris sa signature ; elle ne prétend pas inverser une coordonnée écrasée
sans consulter le registre.

\begin{proposition}[Réalisation autonome du TPK au moyen de transitions explicites]
\label{ix:prop:modelo-tpk-calibrado-ch}
La famille
\[
\begin{aligned}
 \mathfrak T_{\rm TPK}^{\rm cal}:=\bigl(
 &(\widetilde q_{n,r})_{n,r},\pi_{\rm obs},
   \widetilde F_{\rm obs},M_{\rm ph},\\
 &(\mathcal A_{\widetilde q}^{\rm cal})_{\widetilde q},
   (\mathcal E_{\widetilde q}^{\rm cal})_{\widetilde q},
   (\mathcal F_{\widetilde q}^{\rm cal})_{\widetilde q},
   \operatorname{Ext}^{\APP,\rm cal},
   \operatorname{Ext}^{\rm fib,cal},
   \operatorname{Ext}^{\rm enr,cal},\rho^{\rm cal},\\
 &\mathsf{Sel}^{\rm cal},\mathsf{Tra}^{\rm cal},
   \operatorname{Upd}^{\rm cal},\mathsf C_{\mathsf M},
   \kappa_{\mathsf M},\operatorname{gr}_{\mathsf M},
   H_{\rm Carr},Y_{\rm Carr},\\
 &\operatorname{Sig}^{\rm cal},\pi_{\mathcal F}\bigr).
\end{aligned}
\]
est un modèle enrichi du TPK. En particulier, chaque paire du membre de droite de \eqref{ix:eq:transicion-cal-en-tpk-ch} est une arête TPK effective, et non la conséquence inférée d'une condition seulement nécessaire.
\end{proposition}

\begin{proof}
Les étiquettes \((n,r)\) rendent disjointes les copies des porteurs et fixent
sans ambiguïté chaque origine et destination. La relation
\(\operatorname{Ext}^{\APP,\rm cal}\) transporte la cellule complète et
préserve les deux identités euclidiennes ; la relation
\(\operatorname{Ext}^{\rm fib,cal}\) relie les fibres des états
complets antérieur et postérieur : elle transporte feuillet, orientation, préfixe,
cylindre, frontière et registre et incorpore la retenue, la mémoire et la signature
actualisées. La sélection ne modifie pas le registre et le transport concatène
exactement une entrée et produit la présignature de destination
\eqref{ix:eq:prefirma-transporte-cal-ch}. L’actualisation recalcule la signature au moyen de
l’application typée \(\operatorname{Sig}^{\rm cal}_{\widetilde q}\) à partir du
feuillet, de l’orientation, de la retenue, de la route, de la frontière et de la mémoire déjà déterminés. L’application
\(\operatorname{Upd}^{\rm cal}\) est unaire car l’arête contenue dans
\(z_{\varepsilon,r}(x)\) détermine ses cocycles.

Les identités sont les diagonales et les extensions par mots sont les
compositions de \eqref{ix:eq:ext-palabra-calibrada-ch}. La concaténation des
préfixes et des registres est associative ; la mémoire vérifie la loi de cocycle et
la retenue vérifie
\eqref{ix:eq:ley-carry-correcta-ch}. Par conséquent, les deux parenthésages de trois
arêtes produisent le même état. La formule
\eqref{ix:eq:truncamiento-calibrado-definido-ch} définit la troncature sur
les états complets et l’unicité du dernier registre prouve qu’elle retrouve
l’origine dans tous les champs. Ce sont les lois d’identité, de source,
de but, de composition, de mémoire, de retenue et de troncature du TPK enrichi.
En particulier, les relations d’extension sont déterminées comme graphes
d’applications effectives, et non comme conséquences de simples conditions
nécessaires.
\end{proof}

Chaque \(\gamma_{\varepsilon,k_n}\) comprend exactement neuf applications de \(\operatorname{Upd}^{\rm cal}_{e_{\varepsilon,n,r}}\circ
\mathsf{Tra}^{\rm cal}_{e_{\varepsilon,n,r}}\circ
\mathsf{Sel}^{\rm cal}_{e_{\varepsilon,n,r}}\). La survie sera prouvée au moyen d'une queue construite et ne sera pas utilisée pour définir \(G_n\).

\begin{lemma}[Trisection interne effective d'un tour]
\label{ix:lem:elevacion-catalogal-nonadica}
Pour tout \(n\ge0\), tout \(x\in G_n\) et tout \(\varepsilon\in\{-1,0,+1\}\), l'expression \(\gamma_{\varepsilon,k_n}(x)\) est définie par neuf transitions TPK, avec leurs relations synchronisées \(\operatorname{Ext}^{\APP,\rm cal}\) et \(\operatorname{Ext}^{\rm fib,cal}\), et satisfait
\begin{align}
 \rho^{\rm cal}_{k_n+9,k_n}\gamma_{\varepsilon,k_n}(x)&=x,
 \label{ix:eq:truncamiento-injerto-ch}\\
 g\bigl(\gamma_{\varepsilon,k_n}(x)\bigr)
   &=g(x),&
 \jmath_{k_n+9}\!\left(\operatorname{mem}_{k_n+9}
        (\gamma_{\varepsilon,k_n}(x))\right)
   &=\jmath_{k_n}\!\left(\operatorname{mem}_{k_n}(x)\right)+u,
 \label{ix:eq:fase-memoria-injerto-ch}\\
 c\bigl(\gamma_{\varepsilon,k_n}(x)\bigr)
   &=c(x)+\varepsilon\kappa_{\rm Carr}.&&
 \label{ix:eq:carry-injerto-ch}
\end{align}
\label{ix:eq:tres-vueltas-propiedades-ch}
\label{ix:eq:tres-vueltas-acarreo-ch}
Le registre ordonné des neuf arêtes détermine \(\varepsilon\) de manière unique. Les trois images sont donc disjointes. Chaque extrémité appartient à \(G_{n+1}\) et admet à nouveau les trois applications \(\gamma_{-1,k_{n+1}},\gamma_{0,k_{n+1}},
\gamma_{+1,k_{n+1}}\).
\end{lemma}

\begin{proof}
Les formules \eqref{ix:eq:celula-completa-fase-ch} parcourent une fois les neuf phases et renvoient la copie \(a_{n+1,1}\) de \(a(1)\). À chaque étape, la relation APP transporte le sextuplet complet ; \(\mathsf{Sel}^{\rm cal}\) fixe le type TRIT à partir de la marque de phase \(r\) conservée dans l'état ; et \(\mathsf{Tra}^{\rm cal}\) ajoute exactement un symbole au préfixe, un segment au cylindre et sa paire ordonnée d'extrémités à la frontière. Par récurrence sur \(r\), les tuples explicites \(y_{\varepsilon,r}(x)\) et \(z_{\varepsilon,r}(x)\) satisfont les prédicats des relations indiquées, et la mise à jour unaire produit l'état du niveau suivant. L'identité \(\rho^{\rm cal}_{d_{n,r}+1,d_{n,r}}
(x_{\varepsilon,r}(x))=x_{\varepsilon,r-1}(x)\) est la définition \eqref{ix:eq:truncamiento-calibrado-definido-ch} ; son caractère bien défini découle de l'unicité du dernier registre. La composition des neuf applications de troncature récupère \(x\), ce qui prouve \eqref{ix:eq:truncamiento-injerto-ch} sans inverser séparément les champs écrasés.

La phase \(9\to1\) apparaît une seule fois. Par
\eqref{ix:eq:memoria-arista-ch} et la loi de cocycle, la partie linéaire de mémoire
transporte la mémoire préalable et sa partie translationnelle ajoute exactement \(u\)
dans la limite directe. La compatibilité
\(\jmath_{d+1}\jmath_{d+1,d}=\jmath_d\) prouve
\eqref{ix:eq:fase-memoria-injerto-ch} ; le retour concerne la phase, non
l’état enrichi.

Par \eqref{ix:eq:carry-arista-derivado-ch}, huit arêtes ont un incrément nul et
l’arête qui complète \(\mathfrak p_\varepsilon\) possède l’incrément calculé
\(\chi(\varepsilon,\varepsilon)\kappa_{\rm Carr}\). La loi
\eqref{ix:eq:ley-carry-correcta-ch} et
\eqref{ix:eq:carry-no-parametrico-ch} donnent
\eqref{ix:eq:carry-injerto-ch}. Le registre conserve
\((\mathfrak p_\varepsilon,2\varepsilon,r_3(2\varepsilon),
\chi(\varepsilon,\varepsilon))\) ; les trois valeurs de ce registre sont
distinctes même lorsque l’incrément scalaire est nul. C’est pourquoi les images
ne sont pas identifiées et le décodeur de bloc retrouve un unique
\(\varepsilon\).

Les définitions des arêtes dépendent seulement de la nouvelle phase et concatènent,
au lieu de remplacer, préfixe, cylindre, frontière et registre. Elles sont aussi
invariantes sous
\(\jmath_d\operatorname{mem}_d\mapsto
  \jmath_d\operatorname{mem}_d+nu\)
et sous la translation de la retenue
d’entrée. Par conséquent, la même construction s’applique à l’extrémité de
tout tour. Pour chaque état fini, la suite explicite
\[
 x,\quad\gamma_{0,k_n}(x),\quad
 \gamma_{0,k_{n+1}}\gamma_{0,k_n}(x),\quad\ldots
\]
est une queue infinie compatible. Cela démontre la survie de tous les
états construits et la disponibilité réitérée des trois branches, sans
postuler aucune des deux propriétés. Nous pouvons donc, enfin,
identifier, conformément aux définitions de niveau précédentes,
\begin{equation}
 H_{k_n}^{\rm cal}:=G_n
 \label{ix:eq:subarbol-calibrado-superviviente-ch}
\end{equation}
pour tout \(n\ge0\).
\end{proof}

La clôture précédente distingue les trois horloges impliquées. Neuf émissions
ramènent la phase dans \(C_9\) et font avancer la mémoire. Douze tours, soit
\(108\) émissions, ramènent en outre la position du calendrier observable ;
le registre et la mémoire ne sont pas non plus effacés à ce moment. L’opérateur
\(\mathcal K_{\rm ph}\) agit uniquement ensuite, comme reconnaissance de
mémoire réduite des tours déjà construits ; il n’engendre jamais les arêtes de
\(\operatorname{Ext}^{\rm enr,cal}\).

% END OWNER


Soit
\[
 \upsilon:\mathbb F_3\longrightarrow\{-1,0,+1\},\qquad
 \upsilon(0)=0,\quad\upsilon(1)=+1,\quad\upsilon(2)=-1
\]
le représentant équilibré. Pour \(b=(b_1,\ldots,b_6)\in\mathbb F_3^6\), le prolongement de bloc est la composition ordonnée de six tours
\[
 \Gamma_{b,k}^{[54]}:=
 \gamma_{\upsilon(b_6),k+45}\circ\gamma_{\upsilon(b_5),k+36}\circ
 \gamma_{\upsilon(b_4),k+27}\circ\gamma_{\upsilon(b_3),k+18}\circ
 \gamma_{\upsilon(b_2),k+9}\circ\gamma_{\upsilon(b_1),k}.
\tag{39d}
\label{ix:eq:gamma-bloque-54-ch}
\]
Au niveau de bloc, on emploie toujours \(k=k_{6N}=1+54N\) ; chaque facteur a donc le domaine construit par le facteur précédent. Fixons un germe APP canonique \(x_\star\in H^{\rm cal}_1\) et définissons récursivement
\[
 \begin{aligned}
  Y_0&:=\{x_\star\},\qquad Y_N\subseteq H_{1+54N}^{\rm cal},\\
  Y_{N+1}&:=\{\Gamma_{b,1+54N}^{[54]}(x):
                 x\in Y_N,\ b\in\mathbb F_3^6\},\\
  \rho^{\rm blk}_{N+1,N}
    &:=\rho^{\rm cal}_{1+54(N+1),1+54N}
       \!\upharpoonright Y_{N+1},\\
  \operatorname{Ext}^{\rm cal}_N(x)
    &:=\{\Gamma_{b,1+54N}^{[54]}(x):b\in\mathbb F_3^6\}.
 \end{aligned}
\tag{39e}
\label{ix:eq:sincronizacion-54-ch}
\]
La notation \(\rho^{\rm cal}_{b,a}\) désigne la composition de toutes les suppressions
entre les niveaux \(b\) et \(a\). Par construction, chaque prolongement de
\(\operatorname{Ext}^{\rm cal}_N(x)\) est une chaîne de cinquante-quatre
arêtes TPK, et non un saut du quotient de Markov.

Le lecteur est défini directement sur le registre, avant de faire appel à une représentation au moyen de \(\Gamma_b\). Le lemme précédent fournit un décodeur
\[
 \operatorname{dec}_9:\{\text{segments générés de neuf arêtes}\}
 \longrightarrow\{-1,0,+1\},
\]
qui lit le couple TRIT et sa réduction dans le segment. Si
\(\operatorname{Led}^{(9)}_{j,i}(x)\) désigne le \(i\)-ième segment nonadique
du bloc \(j\) du registre de \(x\), on fixe
\begin{equation}
 \beta_{{\rm blk},N}(x):=
 \Bigl(
   \bigl(\upsilon^{-1}\operatorname{dec}_9
          (\operatorname{Led}^{(9)}_{j,i}(x))\bigr)_{i=1}^{6}
 \Bigr)_{j=0}^{N-1}.
 \label{ix:eq:lector-bloques-ch}
\end{equation}
Pour \(N=0\), la valeur est le mot vide. Cette définition ne présuppose pas que la décomposition de \(x\) en composition d'applications \(\Gamma_b\) soit unique.

Spécialisons maintenant la conjugaison arithmétique précédente à
\(p=3\), \(d=6\) et \(A=I_6\). Si
\(\pi_{N+1,N}:(\mathbb Z/3^{N+1}\mathbb Z)^6\to
(\mathbb Z/3^N\mathbb Z)^6\) est la réduction canonique, nous définissons
\begin{equation}
 \Phi_N:=\Psi_N^{-1}\circ\beta_{{\rm blk},N}:
 Y_N\longrightarrow(\mathbb Z/3^N\mathbb Z)^6.
 \label{ix:eq:phi-hensel-bloques-ch}
\end{equation}
Cette définition n’assigne pas une coordonnée après avoir observé une branche : le
lecteur \(\beta_{{\rm blk},N}\) s’obtient à partir du registre TPK déjà construit et
\(\Psi_N^{-1}\) est l’inverse explicite, fixé auparavant, de la tour
arithmétique non filtrée.

\begin{proposition}[Conjugaison finie et recouvrement complet de la fibre]
\label{ix:prop:cobertura-729-seccion}
Pour tout \(N\geq0\),
\[
 \beta_{{\rm blk},N}:Y_N\xrightarrow{\ \cong\ }
 (\mathbb F_3^6)^N
\]
est une bijection compatible avec les troncatures, et pour chaque \(x\in Y_N\)
\[
 \begin{gathered}
 \beta_{{\rm blk},N+1}
 \bigl[\operatorname{Ext}^{\rm cal}_N(x)\bigr]
 =\{\beta_{{\rm blk},N}(x)^\frown b:
       b\in\mathbb F_3^6\},\\
 \operatorname{pref}_N\circ\beta_{{\rm blk},N+1}
 =\beta_{{\rm blk},N}\circ\rho^{\rm blk}_{N+1,N}.
 \end{gathered}
\tag{40$_0$}
\label{ix:eq:cobertura-729-ch}
\]
Les applications \(\Phi_N\) sont des bijections et vérifient l’identité
d’entrelacement
\begin{equation}
 \pi_{N+1,N}\circ\Phi_{N+1}
 =\Phi_N\circ\rho^{\rm blk}_{N+1,N}.
 \label{ix:eq:entrelazamiento-hensel-ext-ch}
\end{equation}
Par conséquent, la tour des histoires engendrées et la tour arithmétique de Hensel sont isomorphes en tant que systèmes projectifs finis.
\end{proposition}

\begin{proof}
L’assertion est immédiate pour \(N=0\). Supposons qu’elle soit vraie au niveau
\(N\). Le lemme précédent construit, pour chaque mot
\(b\in\mathbb F_3^6\), un prolongement
    \(\Gamma_{b,1+54N}^{[54]}(x)\). Si \(b\ne b'\), la première coordonnée par laquelle
ils diffèrent détermine une paire TRIT et un cocycle distincts dans le premier
segment nonadique divergent du registre généalogique. Ce registre conserve
les arêtes ordonnées, et non seulement la somme finale de leurs retenues ; les deux
histoires sont donc distinctes même si les retenues de tours ultérieurs se
compensent. La lecture directe \eqref{ix:eq:lector-bloques-ch} donne alors
\begin{equation}
 \beta_{{\rm blk},N+1}
 \bigl(\Gamma_{b,1+54N}^{[54]}(x)\bigr)
 =\beta_{{\rm blk},N}(x)^\frown b.
 \label{ix:eq:beta-concatenacion-demostrada-ch}
\end{equation}
Les six trisections successives produisent exactement \(3^6\) prolongements distincts. La suppression des cinquante-quatre dernières étapes récupère \(x\), et la suppression du dernier bloc du registre récupère \(\beta_{{\rm blk},N}(x)\). Cela prouve simultanément la bijection et la compatibilité de \(\beta_{{\rm blk},N}\) avec les troncatures.

La proposition~\ref{ix:prop:ch-hensel-bruto} prouve, indépendamment, que
\(\Psi_N\) est bijective et que
\[
 \operatorname{pref}_N\circ\Psi_{N+1}
 =\Psi_N\circ\pi_{N+1,N}.
\]
En composant cette égalité avec la compatibilité de \(\beta_{\rm blk}\) et en
appliquant \(\Psi_N^{-1}\), on obtient
\[
 \pi_{N+1,N}\Psi_{N+1}^{-1}\beta_{{\rm blk},N+1}
 =\Psi_N^{-1}\beta_{{\rm blk},N}\rho^{\rm blk}_{N+1,N},
\]
qui est exactement \eqref{ix:eq:entrelazamiento-hensel-ext-ch}. L’égalité ne
provient pas de la définition d’une coordonnée ad hoc sur chaque branche : elle résulte de la composition
de deux systèmes de bijections construits et compatibles à toutes les
profondeurs.
\end{proof}

La limite projective
\[
 \Omega_\Gamma^{\rm cal}:=
 \varprojlim_N(Y_N,\rho^{\rm blk}_{N+1,N})
\tag{40a$_0$}
\label{ix:eq:omega-gamma-catalogal-ch}
\]
est la limite d'histoires du modèle TPK construit arête par arête dans le lemme précédent. Sa projection par \(\pi_{\rm obs}\) conserve les configurations du calendrier observable, mais l'affirmation ne dépend pas de l'identification de ce modèle autonome à une relation ambiante définie uniquement par des conditions nécessaires. Les bijections finies passent à la limite et produisent l'homéomorphisme cylindrique
\[
 \Lambda:(\mathbb F_3^6)^{\mathbb N}
 \xrightarrow{\ \cong\ }\Omega_\Gamma^{\rm cal},
 \qquad
 \beta_{\rm blk}:=\Lambda^{-1}.
\tag{40b$_0$}
\label{ix:eq:homeomorfismo-calibrado-ch}
\]
La projection canonique est notée
\(\rho^{\rm blk}_{\infty,N}:\Omega_\Gamma^{\rm cal}\to Y_N\).
Ainsi, \(\Lambda\) établit une équivalence compatible avec les
troncatures : chaque ruban ternaire encode une histoire du TPK et chaque
préfixe se prolonge au moyen d’arêtes du même générateur.

Les sous-cylindres sont maintenant définis comme sous-ensembles de la limite, et non comme paires
de suffixes. Pour \(x\in Y_N\) et \(b\in\mathbb F_3^6\), soient
\begin{align}
 \mathscr Z_N(x)
   &:=\{\xi\in\Omega_\Gamma^{\rm cal}:
       \rho^{\rm blk}_{\infty,N}(\xi)=x\},
 \label{ix:eq:cilindro-real-limite-ch}\\
 \mathscr Z_{N+1}(x;b)
   &:=\{\xi\in\Omega_\Gamma^{\rm cal}:
       \rho^{\rm blk}_{\infty,N+1}(\xi)
       =\Gamma_{b,1+54N}^{[54]}(x)\}.
 \label{ix:eq:subcilindro-real-limite-ch}
\end{align}

\begin{theorem}[Partition effective en \(729\) sous-cylindres de raffinement]
\label{ix:thm:particion-729-subcilindros-ch}
Pour tout \(N\ge0\) et tout \(x\in Y_N\), on a l’union disjointe
\begin{equation}
 \mathscr Z_N(x)=
 \bigsqcup_{b\in\mathbb F_3^6}\mathscr Z_{N+1}(x;b).
 \label{ix:eq:particion-real-729-ch}
\end{equation}
Les \(3^6=729\) termes sont des sous-cylindres de raffinement non vides et
ouverts-fermés ; chacun est réalisé par cinquante-quatre arêtes du TPK
construit.
\end{theorem}

\begin{proof}
Soit \(\xi\in\mathscr Z_N(x)\). Sa coordonnée de niveau \(N+1\) appartient à \(Y_{N+1}\), se tronque en \(x\) et, par la bijection \(\beta_{{\rm blk},N+1}\), possède un dernier bloc unique \(b\in\mathbb F_3^6\). L'identité \eqref{ix:eq:beta-concatenacion-demostrada-ch} implique alors \(\rho^{\rm blk}_{\infty,N+1}(\xi)
=\Gamma_{b,1+54N}^{[54]}(x)\) ; par conséquent, \(\xi\) appartient à l'un des termes du membre de droite. L'inclusion inverse est immédiate par troncature. Si deux termes s'intersectaient, leurs coordonnées de niveau \(N+1\) seraient égales et le décodeur du registre donnerait \(b=b'\), de sorte que la réunion est disjointe.

Chaque extrémité \(\Gamma_b^{[54]}(x)\) admet une queue infinie obtenue
en itérant le prolongement de bloc nul ; tous les termes sont donc non
vides. Comme \(Y_{N+1}\) est fini et discret, ce sont des images réciproques de points par
une projection continue de la limite projective et, par conséquent, ils sont ouverts et fermés. La
définition même de \(\Gamma_b^{[54]}\) atteste ses cinquante-quatre arêtes.
\end{proof}

La famille \((\Phi_N)_N\) induit finalement l’homéomorphisme de limites
\begin{equation}
 \Phi_\infty:\Omega_\Gamma^{\rm cal}
 \xrightarrow{\ \cong\ }\mathbb Z_3^6,\qquad
 \Phi_\infty(\xi):=
 \bigl(\Phi_N(\rho^{\rm blk}_{\infty,N}(\xi))\bigr)_{N\ge0},
 \label{ix:eq:phi-infinito-hensel-ext-ch}
\end{equation}
où
\(\mathbb Z_3^6=\varprojlim_N(\mathbb Z/3^N\mathbb Z)^6\).
L’équation \eqref{ix:eq:entrelazamiento-hensel-ext-ch} garantit que le membre
droit est une classe compatible, et les bijections finies prouvent
l’injectivité, la surjectivité et la continuité dans les deux sens. Sous
\(\Phi_\infty\), \(\mathscr Z_N(x)\) est exactement la classe de congruence
\[
 \{z\in\mathbb Z_3^6:z\equiv\Phi_N(x)\pmod{3^N}\},
\]
et les ensembles \(\mathscr Z_{N+1}(x;b)\) sont ses \(729\) classes de raffinement
modulo \(3^{N+1}\). Tel est le raccord effectif entre prolongement TPK,
cylindres de la limite et tour de Hensel.

Pour \(b=(b_1,\ldots,b_6)\in\mathbb F_3^6\), posons
\[
 \nu(b):=\sum_{i=1}^{6}[b_i]_3\,3^{6-i}
 \in\{0,\ldots,728\},
\tag{40c$_0$}
\label{ix:eq:valor-bloque-ch}
\]
où \([b_i]_3\in\{0,1,2\}\) est le représentant positionnel. Les
\(243\) mots visibles du recensement initial sont l’image de la première
émission ; ils ne se confondent pas avec la capacité \(3^6\) de chaque bloc de
prolongement. Enfin,
\[
 X_{\infty,\mathbb Z}^{\rm cal}
 :=\mathbb Z_{\rm HMT}\times\Omega_\Gamma^{\rm cal}.
\]

% END OWNER


Pour un préfixe \(s=(m;b_0,\ldots,b_{N-1})\), définissons dans l’arithmétique
engendrée de \(\mathbb Q_{\rm HMT}\)
\[
 a_N(s):=m+\sum_{q<N}\frac{\nu(b_q)}{729^{q+1}},
 \qquad
 \operatorname{Succ}(s):=
 \{(m;b_0,\ldots,b_{N-1},b):b\in\mathbb F_3^6\}.
\]
Si \((m;\beta_{{\rm blk},N}(x))\) désigne le préfixe obtenu en préfixant
\(m\) au tuple \(\beta_{{\rm blk},N}(x)\),
\eqref{ix:eq:cobertura-729-ch} donne l’égalité typée
\[
 \operatorname{Succ}(m;\beta_{{\rm blk},N}(x))
 =\{(m;\beta_{{\rm blk},N+1}(y)):
       y\in\operatorname{Ext}^{\rm cal}_N(x)\}.
\]
Les cylindres semi-ouverts sont, dès leur première définition, des sous-ensembles de
la complétion HMT :
\[
 I_N^{\rm HMT}(s):=
 \{r\in\mathbb R_{\rm HMT}:
      a_N(s)\le_{\rm HMT}r<_{\rm HMT}a_N(s)+729^{-N}\}.
\]
Ils satisfont matériellement
\[
 I_N^{\rm HMT}(s)=
 \bigsqcup_{y\in\operatorname{Succ}(s)}I_{N+1}^{\rm HMT}(y),
 \qquad |\operatorname{Succ}(s)|=729,
 \qquad \operatorname{diam}_{\rm HMT}I_N^{\rm HMT}(s)
      =729^{-N}\longrightarrow0.
\tag{40a}\label{ix:eq:ch-particion-cilindros}
\]
La convention de frontière n’identifie les deux écritures qu’après
les avoir construites. La partition~\eqref{ix:eq:ch-particion-cilindros}, l’emboîtement et la contraction produisent

\[
\begin{aligned}
 P_{\rm ar}&:X_{\infty,\mathbb Z}^{\rm cal}
     \twoheadrightarrow\mathbb R_{\rm HMT},\\
 P_{\rm ar}(m,\xi)
   &:=\text{l'unique élément de }
     \bigcap_N C_N^{\rm cl}
       \bigl(m;\beta_{{\rm blk},N}(\rho^{\rm blk}_{\infty,N}(\xi))\bigr),\\
 C_N^{\rm cl}(s)&:=\overline{I_N^{\rm HMT}(s)}.
\end{aligned}
\tag{40}\label{ix:eq:ch-proyeccion-arquimediana}
\]

Par conséquent,

\[
 X_{\infty,\mathbb Z}^{\rm cal}/\!\sim_{P_{\rm ar}}
 \;\cong\;\mathbb R_{\rm HMT}.
\tag{41}\label{ix:eq:ch-cociente-recta}
\]

L’autre publication \(P_{\rm exc}\) conserve incidence, frontière et mémoire
vers Paley--Witt--Golay--Leech et, de là, vers
\(V^\natural\)--Monster. Toutes deux proviennent du
même antécédent de~\eqref{ix:eq:ch-limite-inverso} ; ni la coordonnée archimédienne ni la reconnaissance
exceptionnelle ne sélectionnent les germes. Les cinq constructions
\(sp,sol,gau,coh,det\) demeurent consubstantielles au sein de
\(\mathcal C_{\rm cont}^{\rm disc}\) et ne se dissocient pas pour effectuer le passage
cardinal.


% BEGIN OWNER /Users/ruben/Documents/New project/output/REVISION_ARGUMENTAL_APERTURAS_CIERRES_20260915/IX/ES/sections/ch_descenso_por_fibras.tex
\subsubsection{Descente des opérations au quotient d’évaluation}
\label{ix:sec:ch-descenso-fibras}

Le quotient précédent identifie les histoires qui publient la même valeur, et non
leurs registres complets. Le critère suivant caractérise exactement
quand une opération d’histoires induit une opération dans ce quotient sans
choisir une histoire privilégiée dans chaque fibre ; sa formulation
poursuit la construction généalogique établie dans les sections précédentes.

\begin{proposition}[Critère de descente par fibres d’évaluation]
\label{ix:prop:ch-descenso-fibras}
Soit \(q:\Omega\to J\) l’évaluation sur son image d’une carte
généalogique. Soit \(\star:D\to\Omega\) une opération avec
\(D\subseteq\Omega^2\) saturé par les fibres de \(q\times q\).
Il existe une unique opération
\(\overline\star:(q\times q)(D)\to J\) telle que
\[
 q(x\star y)=q(x)\,\overline\star\,q(y)
\]
si et seulement si, pour des paires du domaine,
\[
 q(x)=q(x'),\quad q(y)=q(y')
 \quad\Longrightarrow\quad q(x\star y)=q(x'\star y').
\]
\end{proposition}

\begin{proof}
Si l'opération visible existe, son résultat dépend seulement des deux valeurs, ce qui prouve la nécessité. Pour la suffisance, on définit \(s\,\overline\star\,t=q(x\star y)\), où \(q(x)=s\) et \(q(y)=t\). La saturation conserve le domaine lors du changement de représentants ; la condition énoncée conserve le résultat. La surjectivité de \(q\times q\) sur l'image du domaine prouve l'unicité.
\end{proof}

Une coordonnée ne remplace pas automatiquement une histoire comme argument de
toutes ses opérations : le critère détermine quand cette substitution est
légitime. En particulier, transporter une opération à travers une bijection
ne suffit pas à l’identifier à une opération native antérieure à la lecture.
Si une carte \(\Omega\) est compacte et \(q\) est continue et surjective
sur un espace séparé \(J\), l’évaluation induit en outre un
homéomorphisme \(\Omega/{\sim_q}\to J\) : c’est une bijection continue d’un
compact vers un espace séparé. C’est le \emph{quotient}, et non nécessairement
l’espace des histoires, qui s’identifie topologiquement à \(J\).

% END OWNER


\subsubsection{Reconstruction interne de la carte et relèvement uniforme}

À l'intérieur de \(\mathfrak G_{\rm HMT}(h,m_\infty)\), la reconstruction de \eqref{ix:eq:ch-limite-inverso}--\eqref{ix:eq:ch-cociente-recta} est réalisée avec les mêmes codes et opérateurs. Notons \(\mathbb R_{\rm HMT}^{\mathfrak G}\) la complétion HMT interne. L'unicité du corps ordonné complet archimédien fournit, \textbf{après} l'avoir construit, l'isomorphisme interne
\[
 \eta:\mathbb R_{\rm HMT}^{\mathfrak G}
 \xrightarrow{\ \cong\ }\mathbb R^{\mathfrak G}.
\tag{41$_0$}\label{ix:eq:ch-isomorfismo-interno}
\]
Pour l’intervalle interne correspondant
\(I_N^{\mathfrak G}(s):=
 I_N^{\rm HMT}(s)\cap\mathbb R_{\rm HMT}^{\mathfrak G}\), l’isomorphisme
préserve les extrémités et l’ordre :
\[
 \eta[I_N^{\mathfrak G}(s)]
 =\{r\in\mathbb R^{\mathfrak G}:
      a_N(s)\le r<a_N(s)+729^{-N}\}.
\tag{41$_0$a}\label{ix:eq:ch-intervalos-internos}
\]
La surjectivité de~\eqref{ix:eq:ch-proyeccion-arquimediana} couvre tous les éléments de cette complétion, non
seulement ceux obtenus dans une exécution finie.

Pour que « tous » ne reste pas caché dans la définition du codomaine, on construit
les deux injections qui relient la droite à son ensemble interne des parties. Soit

\[
 E(A)=\sum_{n\ge0}\frac{2\mathbf1_A(n)}{3^{2n+2}}
 \qquad
 \bigl(A\in\mathcal P^{\mathfrak G}(\omega)\bigr).
\tag{41a}\label{ix:eq:ch-incrustacion-cantor}
\]

L’écriture place \(0\) ou \(2\) uniquement aux positions ternaires
paires ; son image est contenue dans \([0,1/4]\). Elle n’atteint donc jamais
l’extrémité droite du cylindre racine et deux sous-ensembles distincts ont une
première position paire distincte : \(E\) est injective et ne subit pas l’ambiguïté de
frontière. Écrivons \(E_{\rm HMT}:=\eta^{-1}\circ E\).

Étant donné \(A\), on part du cylindre racine qui contient \(E_{\rm HMT}(A)\). Si \(s_N(A)\) est le préfixe déjà choisi, écrivons \(s_N(A)=(0;b_0(A),\ldots,b_{N-1}(A))\). La partition~\eqref{ix:eq:ch-particion-cilindros}, transportée par \(\eta^{-1}\), fournit un unique sous-cylindre semi-ouvert \(s_{N+1}(A)\) qui contient \(E_{\rm HMT}(A)\). D'après~\eqref{ix:eq:cobertura-729-ch}, ce prolongement est également un prolongement TPK admissible. Par récurrence,
\[
 s_{N+1}(A)\!\upharpoonright N=s_N(A),\qquad
 E_{\rm HMT}(A)\in I_N^{\rm HMT}(s_N(A))\quad(N<\omega).
\tag{41b$_0$}\label{ix:eq:ch-prefijos-uniformes}
\]
La récurrence est une formule uniforme de paramètres \(A,h\). Comme tous deux sont
internes, le Remplacement dans le modèle du lemme~\ref{ix:thm:ch-6b} forme la suite complète de
blocs \((b_q(A))_{q<\omega}\). Définissons, sans confondre le ruban et la
suite des préfixes,
\[
 \xi_A:=\Lambda\bigl((b_q(A))_{q<\omega}\bigr)
 \in\Omega_\Gamma^{\rm cal,\mathfrak G}.
\]
On obtient ainsi le \textbf{relèvement uniforme du plongement de Cantor}

\[
 Q:\mathcal P^{\mathfrak G}(\omega)
 \longrightarrow X_{\infty,\mathbb Z}^{\rm cal,\mathfrak G},
 \qquad
 Q(A):=(0,\xi_A),\qquad
 P_{\rm ar}(Q(A))=E_{\rm HMT}(A).
\tag{41b}\label{ix:eq:ch-levantamiento-uniforme}
\]

On n’appelle pas \(Q\) l’inverse de \(P_{\rm ar}\) : leur composition publie le
plongement \(E_{\rm HMT}\), non l’identité sur toute la droite. Dans l’autre
direction, soit \((q_n)_{n<\omega}\) l’énumération engendrée de
\(\mathbb Q_{\rm HMT}\) et définissons
\[
 C(r):=\{n<\omega:\eta(q_n)<r\}
 \qquad(r\in\mathbb R^{\mathfrak G}).
\]
La coupure rationnelle appartient à \(\mathcal P^{\mathfrak G}(\omega)\) et \(C\) est injectif par densité. Les injections \(E\) et \(C\), suivies de Cantor--Bernstein interne, prouvent

\[
 P_{\rm ar}
 \bigl[X_{\infty,\mathbb Z}^{\rm cal,\mathfrak G}\bigr]
 =\mathbb R_{\rm HMT}^{\mathfrak G},
 \qquad
 \mathbb R_{\rm HMT}^{\mathfrak G}
 \cong\mathbb R^{\mathfrak G}
 \approx\mathcal P^{\mathfrak G}(\omega).
\tag{41c}\label{ix:eq:ch-recta-potencia}
\]

Ici, un réel arbitraire est utilisé uniquement comme variable universelle pour prouver
la surjectivité d’un arbre déjà construit. Il ne sélectionne ni APP, ni TRIT, ni TPK, ni les
constantes, ni les transitions de cet arbre.

\subsection{Énumération par codes généalogiques et condensation HMT}

La règle~\eqref{ix:eq:ch-operador-def} fournit des noms de dérivation à partir de

\[
 (\ulcorner\varphi\urcorner,
  \sigma(p_0),\ldots,\sigma(p_{k-1})).
\tag{42}\label{ix:eq:ch-codigo-nombres}
\]

Pour tout objet \(a\) de la hiérarchie, soient \(\beta(a)\) sa première
étape et \(\sigma(a)\) son plus petit arbre de dérivation dans l’ordre récursif du
lemme~\ref{ix:thm:ch-6b}. On n’exige pas que \(\sigma(a)\) soit un naturel lorsque l’étape est non
dénombrable. La relation

\[
 a\prec_{\mathfrak G}b
 \quad\Longleftrightarrow\quad
 (\beta(a),\sigma(a))<_{\rm lex}
 (\beta(b),\sigma(b))
\tag{42a}\label{ix:eq:ch-orden-genealogico}
\]

est le bon ordre généalogique global \(<_{\mathfrak G}\) induit par le
constructeur. Les noms répétés sont éliminés en prenant le plus petit ; aucun réel
ni aucune énumération de la droite n’est inséré comme donnée.

\begin{lemma}[Dénombrabilité interne des niveaux antérieurs à \(\omega_1^{\mathfrak G}\)]
\label{ix:thm:ch-7}
Pour tout \(\beta<\omega_1^{\mathfrak G}\), le niveau \(\mathfrak G_\beta(h,m_\infty)\) est dénombrable à l'intérieur de \(\mathfrak G_{\rm HMT}(h,m_\infty)\), et possède une surjection canonique

\[
 e_\beta:\omega\twoheadrightarrow
 \mathfrak G_\beta(h,m_\infty).
\tag{43}\label{ix:eq:ch-enumeracion-niveles}
\]

\end{lemma}
\begin{proof}
Le niveau initial possède un code dénombrable. Si \(e_\beta\) énumère le niveau \(\beta\), les codes de formule et les tuples finis de paramètres de~\eqref{ix:eq:ch-codigo-nombres} forment un ensemble dénombrable et énumèrent le successeur. Si \(\lambda<\omega_1^{\mathfrak G}\) est limite, il existe intérieurement une surjection \(s:\omega\twoheadrightarrow\lambda\) ; on prend la \(<_{\mathfrak G}\)-plus petite, ordre déjà construit dans le lemme~\ref{ix:thm:ch-6b}, et l'on diagonalise \(e_{s(n)}\). Enfin, parmi toutes les surjections ainsi obtenues, on définit \(e_\beta\) comme la \(<_{\mathfrak G}\)-plus petite. La récursion \(\beta\mapsto e_\beta\) est une classe uniformément définissable de \(\mathfrak G\) ; chaque \(e_\beta\), et chaque restriction de la récursion à un ordinal, appartient à \(\mathfrak G\). Elle n'utilise que \(h,u_{h,m}\) et le bon ordre global.
\end{proof}

\begin{lemma}[Condensation des réels internes en dessous de \(\omega_1^{\mathfrak G}\)]
\label{ix:thm:ch-8}
Tout réel de
\(\mathfrak G_{\rm HMT}(h,m_\infty)\) apparaît dans un niveau d’indice inférieur à
\(\omega_1^{\mathfrak G}\).

\end{lemma}
\begin{proof}
Soit \(r\subseteq\omega\) un réel de \(\mathfrak G\) et soit
\(B:=\mathfrak G_0(h,m_\infty)\). Cette base est transitive et dénombrable
intérieurement. Avec les codes des naturels et des paires finies déjà fixés,
\(\operatorname{Ord}\cap B=\omega+1\). Par récurrence,
\[
 \operatorname{Ord}\cap\mathfrak G_\alpha=(\omega+1)+\alpha.
\]
Le successeur ajoute l’ordinal formé par les ordinaux du niveau précédent ;
il ne peut en ajouter un plus grand, car celui-ci ne serait pas un sous-ensemble de ce niveau. Les
limites sont des unions. Nous choisissons \(\Theta\) limite, \(\Theta>\omega^2\),
avec \(r\in\mathfrak G_\Theta\). Ainsi,
\(\operatorname{Ord}\cap\mathfrak G_\Theta=\Theta\), et tout niveau
\(\mathfrak G_\xi\), avec \(\xi<\Theta\), appartient à
\(\mathfrak G_\Theta\). Nous formons intérieurement la structure étendue
\[
 \mathcal A_\Theta=
 \langle\mathfrak G_\Theta,\in,u_{h,m},
         H_\Theta,<_{\mathfrak G}\!\upharpoonright\mathfrak G_\Theta\rangle
\]
avec la relation de niveaux comme prédicat de sa signature,
\(H_\Theta=\{\langle\xi,x\rangle:\xi<\Theta\ \&\
x\in\mathfrak G_\xi\}\), et ses fonctions de Skolem définissables. Soit
\[
 Y=\operatorname{Hull}^{\mathcal A_\Theta}
   (B\cup\{B,r\})\prec\mathcal A_\Theta .
\]
La clôture sous la famille dénombrable de fonctions de Skolem montre à l’intérieur
de \(\mathfrak G\) que \(Y\) est dénombrable. Soit
\(c:Y\rightarrow M\) son effondrement de Mostowski et
\[
 \bar\Theta:=\operatorname{otp}(Y\cap\Theta)
 <\omega_1^{\mathfrak G}.
\]
Comme \(B\subseteq Y\) est transitif, la récurrence sur l’appartenance montre
que \(c\) fixe chaque élément de \(B\) ainsi que \(B\). En particulier,
elle fixe \(\omega,u_{h,m},h,m_\infty\). Puisque
\(r\subseteq\omega\subseteq Y\), on a aussi \(c(r)=r\).

Nous prouvons par récurrence sur \(\xi\in Y\cap\Theta\) que
\[
 c``\bigl(Y\cap\mathfrak G_\xi(h,m_\infty)\bigr)
 =\mathfrak G_{c(\xi)}(h,m_\infty).
\tag{43a}\label{ix:eq:ch-condensacion}
\]
Pour \(\xi\in Y\cap\Theta\), le prédicat \(H_\Theta\) définit
l’ensemble \(\mathfrak G_\xi\) à partir de \(\xi\), de sorte que
\(\mathfrak G_\xi\in Y\). Relativiser chaque formule à cet ensemble donne
\(Y\cap\mathfrak G_\xi\prec\mathfrak G_\xi\), dans le langage dans lequel
le lemme~\ref{ix:thm:ch-6} traduit les définitions HMT.

La base de~\eqref{ix:eq:ch-condensacion} est déjà prouvée car l’effondrement fixe \(B\). À un
successeur, si \(\xi+1\in Y\), on a aussi \(\xi\in Y\) et
\(c(\xi+1)=c(\xi)+1\). Pour
\(a\in Y\cap\mathfrak G_{\xi+1}\), nous fixons une formule finie qui le
définit sur \(\mathfrak G_\xi\). L’affirmation
\[
 a=\{z\in\mathfrak G_\xi:
                  \varphi^{\mathfrak G_\xi}(z,\bar p)\}
\]
est une formule ordinaire sur \(\mathfrak G_\Theta\), avec paramètres \(a,\mathfrak G_\xi\). L'élémentarité fournit les paramètres \(\bar p\) à l'intérieur de \(Y\cap\mathfrak G_\xi\). L'hypothèse de récurrence et l'élémentarité de ce niveau montrent que \(c(a)\) est le sous-ensemble défini par la même formule et les paramètres effondrés sur \(\mathfrak G_{c(\xi)}\). Réciproquement, étant donnés une formule et des paramètres dans le niveau effondré, on relève les paramètres au niveau original. Le sous-ensemble qui y est défini existe, est unique et appartient à \(Y\) par élémentarité. Son effondrement est le sous-ensemble demandé. Cela prouve les deux inclusions au successeur, sans introduire de prédicat uniforme de vérité.

Pour \(\lambda\in Y\cap\Theta\) limite, l’élémentarité fournit,
pour chaque \(x\in Y\cap\mathfrak G_\lambda\), un témoin
\(\xi\in Y\cap\lambda\) avec \(H_\Theta(\xi,x)\). Par conséquent,
\[
 Y\cap\mathfrak G_\lambda
 =\bigcup_{\xi\in Y\cap\lambda}(Y\cap\mathfrak G_\xi),
 \qquad c``(Y\cap\lambda)=c(\lambda).
\]
La continuité de la hiérarchie et l’hypothèse de récurrence donnent~\eqref{ix:eq:ch-condensacion} à la
limite. La cofinalité de \(Y\cap\lambda\) dans l’ordinal
ambiant \(\lambda\) n’est pas requise : ses indices effondrés parcourent \(c(\lambda)\).

Enfin, chaque \(x\in Y\) possède un témoin \(\xi\in Y\cap\Theta\)
avec \(H_\Theta(\xi,x)\), et \(\bar\Theta\) est limite. En prenant les unions,
\[
 M=\bigcup_{\eta<\bar\Theta}\mathfrak G_\eta(h,m_\infty)
  =\mathfrak G_{\bar\Theta}(h,m_\infty).
\]
Comme \(c(r)=r\in M\), on conclut
\(r\in\mathfrak G_{\bar\Theta}\) avec
\(\bar\Theta<\omega_1^{\mathfrak G}\).
\end{proof}

Le lemme ne dit pas qu’un préfixe de longueur finie énumère la droite. Il dit que
chaque objet héréditaire engendré possède un premier niveau transfini et un premier
code généalogique à l’intérieur de ce niveau. C’est exactement l’information que
le réduit extensionnel perd lorsqu’il ne conserve que la valeur publiée.

\subsection{Injection de la droite engendrée dans le premier ordinal non dénombrable de la clôture généalogique}

Soit \(\iota(r)\subseteq\omega\) la coupure rationnelle de
\(r\in\mathbb R_{\rm HMT}^{\mathfrak G}\) par rapport à l’énumération de
\(\mathbb Q_{\rm HMT}\), construite préalablement à partir des entiers et des quotients
APP. L’application \(r\mapsto\iota(r)\) est définissable à partir de \(r\) et de cette énumération,
de sorte que \(\iota(r)\in\mathfrak G\) ; le lemme~\ref{ix:thm:ch-8} la situe à une étape
d’indice dénombrable. Nous définissons

\[
 \begin{aligned}
  \beta(r)&:=\min\{\beta:
       \iota(r)\in\mathfrak G_{\beta+1}(h,m_\infty)\},\\
  n(r)&:=\min\{n:e_{\beta(r)+1}(n)=\iota(r)\},\\
  j_{\rm HMT}(r)&:=\omega\cdot\beta(r)+n(r).
 \end{aligned}
\tag{44}\label{ix:eq:ch-indices-real}
\]

La définition~\eqref{ix:eq:ch-indices-real} est uniforme dans les paramètres \(h,m_\infty\). Comme \(\mathbb R_{\rm HMT}^{\mathfrak G}\) est un ensemble de \(\mathfrak G\), la séparation fixe le domaine de la définition et le remplacement dans \(\mathfrak G\) forme son graphe :
\[
 \{\langle r,\omega\cdot\beta(r)+n(r)\rangle:
 r\in\mathbb R_{\rm HMT}^{\mathfrak G}\}\in\mathfrak G.
\]
Par conséquent, \(j_{\rm HMT}\) est une fonction interne de \(\mathfrak G\).

Par les lemmes~\ref{ix:thm:ch-7} et~\ref{ix:thm:ch-8},
\(j_{\rm HMT}(r)<\omega_1^{\mathfrak G}\) pour tout
réel engendré. Si \(j_{\rm HMT}(r)=j_{\rm HMT}(s)\), la division ordinale par
\(\omega\) récupère le même couple \((\beta,n)\) ; \eqref{ix:eq:ch-enumeracion-niveles} récupère la même coupure
rationnelle et, par complétude, \(r=s\). Par conséquent,

\[
 j_{\rm HMT}:\mathbb R_{\rm HMT}^{\mathfrak G}
 \hookrightarrow\omega_1^{\mathfrak G}
\tag{45}\label{ix:eq:ch-inyeccion-j}
\]

est une injection construite par le rang et le nom de la généalogie HMT. Son existence utilise exactement la clôture sémantique~\eqref{ix:eq:ch-operador-def}, le bon ordre des noms du lemme~\ref{ix:thm:ch-6b} et la condensation~\eqref{ix:eq:ch-condensacion} ; elle ne reçoit pas l'hypothèse du continu pour choisir l'étape. L'identification ensembliste ultérieure décrit cette même machinerie dans une autre notation, sans inverser sa provenance causale.

\subsection{Injection du premier ordinal non dénombrable dans la droite et équivalence cardinale}

Pour chaque \(\beta<\omega_1^{\mathfrak G}\), le lemme~\ref{ix:thm:ch-7} permet de coder dans un
sous-ensemble \(w_\beta\subseteq\omega\) un bon ordre dénombrable de type
\(\beta\). Pour ne pas confondre l’ordre vide et l’ordre à un point, on code
aussi bien le domaine \(D\subseteq\omega\) que sa relation d’ordre strict
\(R\subseteq D^2\). Un appariement bijectif
\(\langle\cdot,\cdot\rangle:\omega^2\to\omega\) étant fixé, le code est
\[
 w(D,R)=\{2n:n\in D\}\cup
        \{2\langle a,b\rangle+1:(a,b)\in R\}.
\]
Les bits pairs récupèrent le domaine et les impairs la relation. On prend le
plus petit réel selon \(<_{\mathfrak G}\) qui encode un tel bon ordre. Comme la propriété
« encode un bon ordre de type \(\beta\) » et le bon ordre global sont internes,
le Remplacement forme la fonction \(\beta\mapsto w_\beta\) à l’intérieur de
\(\mathfrak G\). Les types d’ordre distincts produisent des codes distincts.
L’application ternaire, déjà utilisée dans~\eqref{ix:eq:ch-incrustacion-cantor}, est définie par

\[
 E(w)=\sum_{n\ge0}\frac{2\mathbf1_w(n)}{3^{2n+2}}
\tag{46}\label{ix:eq:ch-incrustacion-ternaria}
\]

et est injective et évite l'extrémité du cylindre racine. Le relèvement~\eqref{ix:eq:ch-levantamiento-uniforme} assure que \(E_{\rm HMT}(w_\beta)\) appartient à la droite publiée par une histoire HMT. Ainsi

\[
 k_{\rm HMT}(\beta):=E_{\rm HMT}(w_\beta),
 \qquad
 k_{\rm HMT}:\omega_1^{\mathfrak G}
 \hookrightarrow\mathbb R_{\rm HMT}^{\mathfrak G}.
\tag{47}\label{ix:eq:ch-inyeccion-k}
\]

Les injections~\eqref{ix:eq:ch-inyeccion-j} et~\eqref{ix:eq:ch-inyeccion-k}, construites sans utiliser l’égalité recherchée,
donnent par Cantor--Bernstein

\[
 \bigl|\mathbb R_{\rm HMT}^{\mathfrak G}\bigr|
 =\aleph_1^{\mathfrak G}.
\tag{48}\label{ix:eq:ch-igualdad-cardinal}
\]

Comme~\eqref{ix:eq:ch-recta-potencia} identifie la droite HMT complète à
\(\mathcal P^{\mathfrak G}(\omega)\), \eqref{ix:eq:ch-igualdad-cardinal} se réécrit exactement sous la forme

\[
 \boxed{
  \mathfrak G_{\rm HMT}(h,m_\infty)
  \models 2^{\aleph_0}=\aleph_1.}
\tag{49}\label{ix:eq:ch-decision}
\]

L'égalité ne provient ni de \(729\), ni d'un produit informel entre cardinalité et ordinalité, ni des chiffres d'une constante, ni d'une fourchette numérique. Elle provient de la couverture de la droite par la limite projective, de la condensation de la clôture généalogique et des deux injections explicites.

\subsection{Quantification sur la totalité des sous-ensembles internes}

Soit
\(A\in\mathcal P^{\mathfrak G}
(\mathbb R_{\rm HMT}^{\mathfrak G})\). Si \(A\) est dénombrable, la
première alternative est satisfaite. S’il ne l’est pas, \(j_{\rm HMT}[A]\) est un sous-ensemble non
dénombrable de \(\omega_1^{\mathfrak G}\). Il ne peut être borné par un ordinal
dénombrable, puisque chaque segment initial de \(\omega_1\) est dénombrable. Son
énumération croissante est de type \(\omega_1^{\mathfrak G}\), et l’inverse de
\(j_{\rm HMT}\) sur son image produit une injection
\(\omega_1^{\mathfrak G}\hookrightarrow A\). L’inclusion
\(A\subseteq\mathbb R_{\rm HMT}^{\mathfrak G}\), \eqref{ix:eq:ch-igualdad-cardinal} et
Cantor--Bernstein donnent

\[
 \boxed{
 \forall A\in\mathcal P^{\mathfrak G}
   (\mathbb R_{\rm HMT}^{\mathfrak G}),\qquad
 |A|\le\aleph_0
 \quad\text{ou}\quad
 |A|=\bigl|\mathbb R_{\rm HMT}^{\mathfrak G}\bigr|.}
\tag{50}\label{ix:eq:ch-dicotomia}
\]

On quantifie ici la puissance complète du continu intégral engendré par
HMT. Le transport bijectif du pliage avec mémoire transporte~\eqref{ix:eq:ch-dicotomia} dans n’importe laquelle
de ses représentations sans perdre de parties ; il ne remplace pas les injections
\eqref{ix:eq:ch-inyeccion-j}--\eqref{ix:eq:ch-inyeccion-k}.

\subsection{Représentation de la clôture généalogique comme \texorpdfstring{\(L[h,m_\infty]\)}{L[h,m-infinito]}}

Ce n'est qu'après~\eqref{ix:eq:ch-decision}--\eqref{ix:eq:ch-dicotomia} que l'on attribue à~\eqref{ix:eq:ch-clausura-transfinita} son nom dans la notation ensembliste usuelle. Soit \(L[u_{h,m}]\) la hiérarchie constructible relative au réel déjà produit dans~\eqref{ix:eq:ch-codigo-conjunto}. Comme la signature HMT est dénombrable et que tous ses opérateurs sont des relations de graphe codées par \(h\), le lemme~\ref{ix:thm:ch-6} élimine leurs symboles uniformément. On obtient l'égalité de classes

% Se separa el contador hipervinculado de la etiqueta editorial (51) para
% evitar dos destinos PDF homónimos en la última ecuación numerada del capítulo.
\stepcounter{equation}
\[
 \mathfrak G_{\rm HMT}(h,m_\infty)
 =L[u_{h,m}]=L[h,m_\infty].
\tag{51}\label{ix:eq:ch-representacion-L}
\]

Pour prouver la première inclusion, on renforce la récurrence : pour chaque
ordinal \(\beta\), tant \(\mathfrak G_\beta\) que la séquence
\(\langle\mathfrak G_\xi:\xi\leq\beta\rangle\) appartiennent à
\(L[u_{h,m}]\). La base est la clôture transitive d’un ensemble construit
à partir de \(\omega\) et de \(u_{h,m}\) ; les paramètres \(h,m_\infty\) se récupèrent
à partir de \(u_{h,m}\). Au successeur, le niveau précédent est un \emph{ensemble
interne}, non seulement une sous-classe de \(L[u_{h,m}]\). Sa relation de satisfaction
et l’élimination uniforme~\eqref{ix:eq:ch-eliminacion-uniforme} permettent de former intérieurement l’ensemble de
ses sous-ensembles définissables. À une limite, la récursion transfinie,
le Remplacement et l’Union de \(L[u_{h,m}]\) forment la séquence précédente et son union.
La transitivité donne alors \(\mathfrak G\subseteq L[u_{h,m}]\).

Réciproquement, le lemme~\ref{ix:thm:ch-6b} fournit un modèle interne
transitif de ZF, contenant tous les ordinaux et \(u_{h,m}\) comme élément.
La récurrence qui construit \(L[u_{h,m}]\) s’effectue dans ce modèle : à chaque
étape, la satisfaction sur le niveau précédent et la définissabilité relative sont
absolues entre les deux domaines transitifs. Par récurrence, ses niveaux
coïncident avec ceux de la construction ambiante et appartiennent à
\(\mathfrak G\). Ainsi, \(L[u_{h,m}]\subseteq\mathfrak G\). Enfin,
\eqref{ix:eq:ch-codigo-conjunto} rend \(u_{h,m}\) et la paire \((h,m_\infty)\)
interdéfinissables de manière récursive primitive, d’où découle la seconde
égalité.

L'équation~\eqref{ix:eq:ch-representacion-L} identifie explicitement la clôture à la hiérarchie constructible relative au paramètre engendré. Cette représentation ne modifie pas la provenance causale : \(L[u]\) ne sélectionne pas le registre généalogique et ne définit ni les prolongements ni la droite ; la décision de l'hypothèse du continu résulte de la règle d'exhaustivité et des applications démontrées entre \eqref{ix:eq:ch-clausura-transfinita} et~\eqref{ix:eq:ch-dicotomia}. La composition \eqref{ix:eq:ch-codigo-conjunto}, \eqref{ix:eq:ch-clausura-transfinita}, \eqref{ix:eq:ch-equivalencia-semantica} et \eqref{ix:eq:ch-decision}--\eqref{ix:eq:ch-representacion-L}, avec le théorème~\ref{ix:thm:decision-continuo-hmt}, offre une expression complète et uniforme de l'univers réalisé, de la coïncidence sémantique et de la clôture cardinale ; aucune de ces représentations ne remplace la génération APP--TRIT--TPK.

\subsection{Théorème de résolution affirmative de l’hypothèse du continu}

\begin{theorem}[Décision du continu dans HMT]
\label{ix:thm:decision-continuo-hmt}
Pour toute histoire réalisée
\(m_\infty\) produite par l’action conjointe
\(\mathrm{APP}\to\mathrm{TRIT}\to\mathrm{TPK}\to
X_\infty^{\rm enr}\to\mathcal C_{\rm cont}^{\rm disc}\), l’extension
généalogique HMT~\eqref{ix:eq:ch-clausura-transfinita} engendre la reconstruction interne de la droite intégrale
\eqref{ix:eq:ch-proyeccion-arquimediana}, \(\mathbb R_{\rm HMT}^{\mathfrak G}\), construit les injections
\eqref{ix:eq:ch-inyeccion-j} et~\eqref{ix:eq:ch-inyeccion-k}, et satisfait
\(2^{\aleph_0}=\aleph_1\) sur toute sa puissance interne. HMT
démontre donc affirmativement l’hypothèse du continu pour le continu
intégral qu’elle définit et engendre. La notation \(L[h,m_\infty]\) est une
représentation ultérieure de cette hiérarchie, non la prémisse qui décide du
résultat.
\end{theorem}


\subsection{Critères de réfutation de la clôture cardinale}

La liste suivante identifie les défauts qui invalideraient une étape de la clôture
précédente. Les contrôles de construction, de sémantique et de
cardinalité sont distingués :

\begin{enumerate}
\item il existe un niveau \(N\) ou un état survivant sans successeur compatible dans \eqref{ix:eq:ch-extension-no-vacia} ;
\item \(u_{h,m}\) ne code pas héréditairement \(h,m_\infty\), ou une relation
de graphe de \(\Sigma_{\rm HMT}\) n’admet pas l’élimination uniforme~\eqref{ix:eq:ch-eliminacion-uniforme} ;
\item un objet admis par les règles sémantiques HMT ne possède pas de terme dans \eqref{ix:eq:ch-terminos}, ou un terme de~\eqref{ix:eq:ch-terminos} publie un objet étranger à ces règles, rompant \eqref{ix:eq:ch-equivalencia-semantica} ;
\item \(\mathfrak G\) ne satisfait pas l’un des axiomes utilisés, son bon
ordre définissable échoue ou une supposée énumération \(e_\beta\) n’est pas interne et
uniforme ;
\item les sous-cylindres de~\eqref{ix:eq:ch-particion-cilindros} ne forment pas une partition disjointe exhaustive, leurs cylindres cessent de se contracter ou la récursion~\eqref{ix:eq:ch-prefijos-uniformes} ne produit pas d'histoire compatible interne ;
\item l’effondrement du lemme~\ref{ix:thm:ch-8} ne commute pas avec une étape successeur ou limite, ou il existe
un réel interne qui n’apparaît à aucun niveau d’indice inférieur à
\(\omega_1^{\mathfrak G}\) ;
\item la sélection \(\beta\mapsto w_\beta\) n’est pas interne, uniforme ou conserve
le type d’ordre ;
\item deux réels différents ont le même couple généalogique \((\beta,n)\) dans
\eqref{ix:eq:ch-indices-real}, deux ordinaux différents produisent le même code dans~\eqref{ix:eq:ch-inyeccion-k}, ou
l’une des injections \(E,C,j_{\rm HMT},k_{\rm HMT}\) échoue ;
\item il existe une partie interne non dénombrable de la droite dont le rang sous \(j_{\rm HMT}\) est borné dans \(\omega_1^{\mathfrak G}\) ;
\item une étape générative antérieure au recouvrement~\eqref{ix:eq:ch-particion-cilindros} utilise l’hypothèse du continu,
\(\mathbb R\), une constante cible, le Comité des données du Conseil
international des sciences (CODATA), le Particle Data Group (PDG) ou une identification physique
ultérieure comme antécédent causal pour définir un germe, une relation
\(\mathscr L_\tau\), \(\operatorname{Ext}\), un coefficient, l’admissibilité
ou un prolongement de l’arbre. Le relèvement
ultérieur~\eqref{ix:eq:ch-prefijos-uniformes} ne constitue pas une telle inversion : un réel arbitraire déjà construit y est une variable
universelle et sélectionne, dans la partition exhaustive préalablement
démontrée, l’unique sous-cylindre qui le contient.
\end{enumerate}

Aucune exécution limitée à un nombre prédéterminé de niveaux ne remplace le quantificateur universel du système projectif ; une rupture effective de compatibilité à une certaine profondeur le réfuterait. Aucune notation ensembliste ultérieure ne prouve à elle seule~\eqref{ix:eq:ch-decision}. La clôture utilise conjointement le lemme~\ref{ix:thm:ch-6b}, l'identification droite--ensemble des parties de~\eqref{ix:eq:ch-recta-potencia}, les lemmes~\ref{ix:thm:ch-7}--\ref{ix:thm:ch-8} et les deux injections construites à partir de la généalogie.
